\documentclass[11pt,fleqn]{book} % Default font size and left-justified equations


%-=-=-=-=-=-=-=-=-=-=-=-=-=-=-=-=-=-=-=-=-=-=-=-=
%	PACKAGES
%-=-=-=-=-=-=-=-=-=-=-=-=-=-=-=-=-=-=-=-=-=-=-=-=

\usepackage[top=3cm,bottom=3cm,left=3cm,right=3cm,headsep=10pt,a4paper]{geometry} % Page margins

\usepackage{graphicx} % Required for including pictures

\usepackage{lipsum} % Inserts dummy text
\usepackage{subfig}
\usepackage{tikz} % Required for drawing custom shapes

%\usepackage[english]{babel} % English language/hyphenation

\usepackage{enumitem} % Customize lists
\setlist{nolistsep} % Reduce spacing between bullet points and numbered lists

\usepackage{booktabs} % Required for nicer horizontal rules in tables

%\usepackage{sagetex}

\usepackage{subfiles} % Required for use with SageMath Cloud

\usepackage{xcolor} % Required for specifying colors by name
\definecolor{ocre}{RGB}{0,126,196} % Define the orange color used for highlighting throughout the book

%-=-=-=-=-=-=-=-=-=-=-=-=-=-=-=-=-=-=-=-=-=-=-=-=
%	MATHEMATICS SETTINGS
%-=-=-=-=-=-=-=-=-=-=-=-=-=-=-=-=-=-=-=-=-=-=-=-=

\usepackage{
amsmath,
amssymb,
amsfonts,
amsthm,
cancel,
colortbl,
longtable,
pdflscape,
setspace,
mathtools,
pgfplots
}

%-=-=-=-=-=-=-=-=-=-=-=-=-=-=-=-=-=-=-=-=-=-=-=-=
%	STOCKHOLMS COLOR THEME 20151012-071009
%-=-=-=-=-=-=-=-=-=-=-=-=-=-=-=-=-=-=-=-=-=-=-=-=

%== BLUE
	\definecolor{sthlmLightBlue}{RGB}{214,237,252} % HEX #d6edfc
		\newcommand{\csthlmLightBlue}[1]{{\color{sthlmLightBlue}{#1}}}
	\definecolor{sthlmBlue}{RGB}{0,110,191} % HEX #006ebf
		\newcommand{\csthlmBlue}[1]{{\color{sthlmBlue}{#1}}}
	%\definecolor{sthlmDarkBlue}{RGB}{} %HEX #
		%\newcommand{\csthlmDarkBlue}[1]{{\color{sthlmDarkBlue}{#1}}}

%== GREEN

	\definecolor{sthlmLightGreen}{RGB}{213,247,244} % HEX #0d5f7f4
		\newcommand{\csthlmLightGreen}[1]{{\color{sthlmLightGreen}{#1}}}
	\definecolor{sthlmGreen}{RGB}{0,134,127} % #00867f
		\newcommand{\csthlmGreen}[1]{{\color{sthlmGreen}{#1}}}

%== GREY
	%\definecolor{sthlmLightGrey}{RGB}{}
		%\newcommand{\csthlmLightGrey}[1]{{\color{sthlmLightGrey}{#1}}}
	\definecolor{sthlmGrey}{RGB}{245,243,238} % HEX #f5f3ee
		\newcommand{\csthlmGrey}[1]{{\color{sthlmGrey}{#1}}}
	\definecolor{sthlmDarkGrey}{RGB}{51,51,51} % HEX #333333
		\newcommand{\csthlmDarkGrey}[1]{{\color{sthlmDarkGrey}{#1}}}

%== ORANGE
	\definecolor{sthlmLightOrange}{RGB}{255,215,210} % HEX #ffd7d2
		\newcommand{\csthlmLightOrange}[1]{{\color{sthlmLightOrange}{#1}}}
	\definecolor{sthlmOrange}{RGB}{221,74,44} % HEX #dd4a2c
		\newcommand{\csthlmOrange}[1]{{\color{sthlmOrange}{#1}}}

%== PURPLE

	\definecolor{sthlmLightPurple}{RGB}{241,230,252} % HEX #f1e6fc
		\newcommand{\csthlmLightPurple}[1]{{\color{sthlmLightPurple}{#1}}}
	\definecolor{sthlmPurple}{RGB}{93,35,125} % HEX #5d237d
		\newcommand{\csthlmPurple}[1]{{\color{sthlmPurple}{#1}}}

%== RED

	\definecolor{sthlmLightRed}{RGB}{254,222,237} % HEX #c40064
		\newcommand{\csthlmLightRed}[1]{{\color{sthlmLightRed}{#1}}}
	\definecolor{sthlmRed}{RGB}{196,0,100} % HEX #fedeed
		\newcommand{\csthlmRed}[1]{{\color{sthlmRed}{#1}}}
	%\definecolor{sthlmDarkRed}{RGB}{} % HEX #
		%\newcommand{\csthlmDarkRed}[1]{{\color{sthlmDarkRed}{#1}}}

%== YELLOW

	%\definecolor{sthlmLightYellow}{RGB}{}
		%\newcommand{\csthlmLightYellow}[1]{{\color{sthlmLightYellow}{#1}}}
	\definecolor{sthlmYellow}{RGB}{252,191,10} % HEX #fcbf0a
		\newcommand{\csthlmYellow}[1]{{\color{sthlmYellow}{#1}}}

%-=-=-=-=-=-=-=-=-=-=-=-=-=-=-=-=-=-=-=-=-=-=-=-=
%	ISSR COLORS
%-=-=-=-=-=-=-=-=-=-=-=-=-=-=-=-=-=-=-=-=-=-=-=-=

%== BLUE
	\definecolor{issrBlue}{RGB}{0,111,174} % HEX #0066cc
		\newcommand{\cissrBlue}[1]{{\color{issrBlue}{#1}}}

%== GREY
	\definecolor{issrGrey}{RGB}{167,169,172} % HEX #999999
		\newcommand{\cissrGrey}[1]{{\color{issrGrey}{#1}}}

%-=-=-=-=-=-=-=-=-=-=-=-=-=-=-=-=-=-=-=-=-=-=-=-=
%	GLOBAL COLOR THEME
%-=-=-=-=-=-=-=-=-=-=-=-=-=-=-=-=-=-=-=-=-=-=-=-=

%== BLUE

	\newcommand{\cLightBlue}[1]{{\color{sthlmLightBlue}{#1}}}
	\newcommand{\crLightBlue}{\color{sthlmLightBlue}}
	\newcommand{\cnLightBlue}{sthlmLightBlue}
	\newcommand{\cBlue}[1]{{\color{sthlmBlue}{#1}}}
	\newcommand{\crBlue}{\color{sthlmBlue}}
	\newcommand{\cnBlue}{sthlmBlue}
	%\newcommand{\cDarkBlue}[1]{{\color{}{#1}}}
	%\newcommand{\crDarkBlue}{\color{}}
	%\newcommand{\cnDarkBlue}{}

%== GREEN

	%\newcommand{\cLightGreen}[1]{{\color{}{#1}}}
	%\newcommand{\crLightGreen}{\color{}}
	\newcommand{\cnLightGreen}{sthlmLightGreen}
	\newcommand{\cGreen}[1]{{\color{sthlmGreen}{#1}}}
	\newcommand{\crGreen}{\color{sthlmGreen}}
	\newcommand{\cnGreen}{sthlmGreen}
	%\newcommand{\cDarkGreen}[1]{{\color{}{#1}}}
	%\newcommand{\crDarkGreen}{\color{}}
	%\newcommand{\cnDarkGreen}{sthlmDarkGreen}

%== GREY

	\newcommand{\cLightGrey}[1]{{\color{sthlmLightGrey}{#1}}}
	\newcommand{\crLightGrey}{\color{sthlmLightGrey}}
	\newcommand{\cnLightGrey}{sthlmLightGrey}
	\newcommand{\cGrey}[1]{{\color{sthlmGrey}{#1}}}
	\newcommand{\crGrey}{\color{sthlmGrey}}
	\newcommand{\cnGrey}{sthlmGrey}
	\newcommand{\cDarkGrey}[1]{{\color{sthlmDarkGrey}{#1}}}
	\newcommand{\crDarkGrey}{\color{sthlmDarkGrey}}
	\newcommand{\cnDarkGrey}{sthlmDarkGrey}

%== ORANGE

	\newcommand{\cLightOrange}[1]{{\color{sthlmLightOrange}{#1}}}
	\newcommand{\crLightOrange}{\color{sthlmLightOrange}}
	\newcommand{\cnLightOrange}{sthlmLightOrange}
	\newcommand{\cOrange}[1]{{\color{sthlmOrange}{#1}}}
	\newcommand{\crOrange}{\color{sthlmOrange}}
	\newcommand{\cnOrange}{sthlmOrange}
	%\newcommand{\cDarkOrange}[1]{{\color{}{#1}}}
	%\newcommand{\crDarkOrange}{\color{}}
	%\newcommand{\cnDarkOrange}{}

%== PURPLE

	\newcommand{\cLightPurple}[1]{{\color{sthlmLightPurple}{#1}}}
	\newcommand{\crLightPurple}{\color{sthlmLightPurple}}
	\newcommand{\cnLightPurple}{sthlmLightPurple}
	\newcommand{\cPurple}[1]{{\color{sthlmPurple}{#1}}}
	\newcommand{\crPurple}{\color{sthlmPurple}}
	\newcommand{\cnPurple}{sthlmPurple}
	%\newcommand{\cDarkPurple}[1]{{\color{}{#1}}}
	%\newcommand{\crDarkPurple}{\color{}}
	%\newcommand{\cnDarkPurple}{}

%== RED

	\newcommand{\cLightRed}[1]{{\color{sthlmLightRed}{#1}}}
	\newcommand{\crLightRed}{\color{sthlmLightRed}}
	\newcommand{\cnLightRed}{sthlmLightRed}
	\newcommand{\cRed}[1]{{\color{sthlmRed}{#1}}}
	\newcommand{\crRed}{\color{sthlmRed}}
	\newcommand{\cnRed}{sthlmRed}
	\newcommand{\cDarkRed}[1]{{\color{sthlmDarkRed}{#1}}}
	\newcommand{\crDarkRed}{\color{sthlmDarkRed}}
	\newcommand{\cnDarkRed}{sthlmDarkRed}

%== YELLOW

	%\newcommand{\cLightYellow}[1]{{\color{}{#1}}}
	%\newcommand{\crLightYellow}{\color{}}
	%\newcommand{\cnLightYellow}{sthlmLightYellow}
	\newcommand{\cYellow}[1]{{\color{sthlmYellow}{#1}}}
	\newcommand{\crYellow}{\color{sthlmYellow}}
	\newcommand{\cnYellow}{sthlmYellow}
	%\newcommand{\cDarkYellow}[1]{{\color{}{#1}}}
	%\newcommand{\crDarkYellow}{\color{}}
	%\newcommand{\cnDarkYellow}{}

%-=-=-=-=-=-=-=-=-=-=-=-=-=-=-=-=-=-=-=-=-=-=-=-=
%	GLOBAL UNITS
%-=-=-=-=-=-=-=-=-=-=-=-=-=-=-=-=-=-=-=-=-=-=-=-=

\newcommand{\unit}[1]{\ensuremath{\, \mathrm{#1}}}
\newcommand{\doneit}{\textcolor{sthlmRed}{$\odot$}}

%-=-=-=-=-=-=-=-=-=-=-=-=-=-=-=-=-=-=-=-=-=-=-=-=
%	GLOBAL MATH 20151012-065906
%-=-=-=-=-=-=-=-=-=-=-=-=-=-=-=-=-=-=-=-=-=-=-=-=

\newcommand{\limit}{\displaystyle\lim}
\newcommand{\impart}{\textrm{Im}}
\newcommand{\repart}{\textrm{Re}}
\newcommand{\naturalnum}{\mathbb{N}}
\newcommand{\integernum}{\mathbb{Z}}
\newcommand{\rationalnum}{\mathbb{Q}}
\newcommand{\realnum}{\mathbb{R}}
\newcommand{\complexnum}{\mathbb{C}}
\newcommand{\degree}{\ensuremath{{}^{\circ}}\xspace}
%\newcommand{\degree}{^{\circ}} Depricated by above
\newcommand{\rad}{\unit{rad}}
\newcommand{\dd}{\mathrm{d}}
\newcommand{\dx}{\mathrm{d}x}
\newcommand{\dy}{\mathrm{d}y}
\newcommand{\dr}{\mathrm{d}r}
\newcommand{\dtheta}{\mathrm{d}\theta}
\newcommand{\dydx}{\dfrac{\mathrm{d}y}{\mathrm{d}x}}
\newcommand{\dydt}{\dfrac{\mathrm{d}y}{\mathrm{d}t}}
\newcommand{\dydu}{\dfrac{\mathrm{d}y}{\mathrm{d}u}}
\newcommand{\dudx}{\dfrac{\mathrm{d}u}{\mathrm{d}x}}
\newcommand{\dxdt}{\dfrac{\mathrm{d}x}{\mathrm{d}t}}
\newcommand{\decf}{\cRed{$\searrow$}}
\newcommand{\incf}{\cGreen{$\nearrow$}}
\newcommand{\stationary}{$\longrightarrow$}
\newcommand{\vect}[1]{\boldsymbol{#1}}
\newcommand{\farg}[1]{{\color{sthlmBlue}{\left[\boldsymbol{#1}\right]}}}
\newcommand{\set}[1]{\left\{#1\right\}}
\newcommand{\sname}[1]{\texttt{#1}}
\newcommand{\inv}{^{-1}}
\newcommand{\cb}{$\Box$}
\newcommand{\suchthat}{\;\ifnum\currentgrouptype=16 \middle\fi|\;}
\newcommand{\chide}[1]{{\color{sthlmDarkGrey}{#1}}}
\newcommand{\lesssteps}{\cGreen{\textbf{Less Steps Solution:}}}
\newcommand{\solnsteps}{\cGreen{\textbf{Solution:}}}
\newcommand{\qdependlist}{\cDarkGrey{\textbf{Dependencies:}}}
\newcommand{\ihat}{\hat{\imath}}
\newcommand{\jhat}{\hat{\jmath}}
\newcommand{\khat}{\hat{k}}
\DeclareMathOperator{\lcm}{lcm}

\newcommand{\alert}[1]{{\color{sthlmRed}{#1}}}
\newcommand{\mathref}[1]{{\color{sthlmPurple}{#1}}}
\newcommand{\cellyes}{\cellcolor{sthlmGreen!30}Yes}
\newcommand{\cellno}{\cellcolor{sthlmRed!30}No}

\usepackage{multicol}

%-=-=-=-=-=-=-=-=-=-=-=-=-=-=-=-=-=-=-=-=-=-=-=-=
%	FONTS UNITS
%-=-=-=-=-=-=-=-=-=-=-=-=-=-=-=-=-=-=-=-=-=-=-=-=

%\usepackage{avant} % Use the Avantgarde font for headings
%\usepackage{times} % Use the Times font for headings

\usepackage[full]{textcomp}
\usepackage{newpxtext} % osf for text, not math
\usepackage{cabin} % sans serif
\usepackage[varqu,varl]{inconsolata} % sans serif typewriter
\usepackage[bigdelims,vvarbb]{newpxmath} % bb from STIX
\usepackage[cal=boondoxo]{mathalfa} % mathcal

\usepackage{microtype} % Slightly tweak font spacing for aesthetics
\usepackage[utf8]{inputenc} % Required for including letters with accents
\usepackage[T1]{fontenc} % Use 8-bit encoding that has 256 glyphs

%-=-=-=-=-=-=-=-=-=-=-=-=-=-=-=-=-=-=-=-=-=-=-=-=
%	BIBLIOGRAPHY AND INDEX
%-=-=-=-=-=-=-=-=-=-=-=-=-=-=-=-=-=-=-=-=-=-=-=-=
\usepackage[autostyle]{csquotes}
\usepackage[backend=biber,natbib=true,style=alphabetic,citestyle=numeric,sorting=nyt,sortcites=true,autopunct=true,hyperref=true,abbreviate=false,backref=true]{biblatex}
\addbibresource{bibliography.bib} % BibTeX bibliography file
\defbibheading{bibempty}{}

\usepackage{calc} % For simpler calculation - used for spacing the index letter headings correctly
\usepackage{makeidx} % Required to make an index
\makeindex % Tells LaTeX to create the files required for indexing

%-=-=-=-=-=-=-=-=-=-=-=-=-=-=-=-=-=-=-=-=-=-=-=-=
%	MAIN TABLE OF CONTENTS
%-=-=-=-=-=-=-=-=-=-=-=-=-=-=-=-=-=-=-=-=-=-=-=-=

\usepackage{titletoc} % Required for manipulating the table of contents

\contentsmargin{0cm} % Removes the default margin

% Part text styling
\titlecontents{part}[0cm]
{\addvspace{20pt}\centering\large\bfseries}
{}
{}
{}

% Chapter text styling
\titlecontents{chapter}[1.25cm] % Indentation
{\addvspace{12pt}\large\sffamily\bfseries} % Spacing and font options for chapters
{\color{ocre!60}\contentslabel[\Large\thecontentslabel]{1.25cm}\color{ocre}} % Chapter number
{\color{ocre}}
{\color{ocre!60}\normalsize\;\titlerule*[.5pc]{.}\;\thecontentspage} % Page number

% Section text styling
\titlecontents{section}[1.25cm] % Indentation
{\addvspace{3pt}\sffamily\bfseries} % Spacing and font options for sections
{\contentslabel[\thecontentslabel]{1.25cm}} % Section number
{}
{\hfill\color{black}\thecontentspage} % Page number
[]

% Subsection text styling
\titlecontents{subsection}[1.25cm] % Indentation
{\addvspace{1pt}\sffamily\small} % Spacing and font options for subsections
{\contentslabel[\thecontentslabel]{1.25cm}} % Subsection number
{}
{\ \titlerule*[.5pc]{.}\;\thecontentspage} % Page number
[]

% List of figures
\titlecontents{figure}[0em]
{\addvspace{-5pt}\sffamily}
{\thecontentslabel\hspace*{1em}}
{}
{\ \titlerule*[.5pc]{.}\;\thecontentspage}
[]

% List of tables
\titlecontents{table}[0em]
{\addvspace{-5pt}\sffamily}
{\thecontentslabel\hspace*{1em}}
{}
{\ \titlerule*[.5pc]{.}\;\thecontentspage}
[]

%----------------------------------------------------------------------------------------
%	MINI TABLE OF CONTENTS IN PART HEADS
%----------------------------------------------------------------------------------------

% Chapter text styling
\titlecontents{lchapter}[0em] % Indenting
{\addvspace{15pt}\large\sffamily\bfseries} % Spacing and font options for chapters
{\color{ocre}\contentslabel[\Large\thecontentslabel]{1.25cm}\color{ocre}} % Chapter number
{}
{\color{ocre}\normalsize\sffamily\bfseries\;\titlerule*[.5pc]{.}\;\thecontentspage} % Page number

% Section text styling
\titlecontents{lsection}[0em] % Indenting
{\sffamily\small} % Spacing and font options for sections
{\contentslabel[\thecontentslabel]{1.25cm}} % Section number
{}
{}

% Subsection text styling
\titlecontents{lsubsection}[.5em] % Indentation
{\normalfont\footnotesize\sffamily} % Font settings
{}
{}
{}

%-=-=-=-=-=-=-=-=-=-=-=-=-=-=-=-=-=-=-=-=-=-=-=-=
%	PAGE HEADINGS
%-=-=-=-=-=-=-=-=-=-=-=-=-=-=-=-=-=-=-=-=-=-=-=-=

\usepackage{fancyhdr} % Required for header and footer configuration

\pagestyle{fancy}
\renewcommand{\chaptermark}[1]{\markboth{\sffamily\normalsize\bfseries\chaptername\ \thechapter.\ #1}{}} % Chapter text font settings
\renewcommand{\sectionmark}[1]{\markright{\sffamily\normalsize\thesection\hspace{5pt}#1}{}} % Section text font settings
\fancyhf{} \fancyhead[LE,RO]{\sffamily\normalsize\thepage} % Font setting for the page number in the header
\fancyhead[LO]{\rightmark} % Print the nearest section name on the left side of odd pages
\fancyhead[RE]{\leftmark} % Print the current chapter name on the right side of even pages
\renewcommand{\headrulewidth}{0.5pt} % Width of the rule under the header
\addtolength{\headheight}{2.5pt} % Increase the spacing around the header slightly
\renewcommand{\footrulewidth}{0pt} % Removes the rule in the footer
\fancypagestyle{plain}{\fancyhead{}\renewcommand{\headrulewidth}{0pt}} % Style for when a plain pagestyle is specified

% Removes the header from odd empty pages at the end of chapters
\makeatletter
\renewcommand{\cleardoublepage}{
\clearpage\ifodd\c@page\else
\hbox{}
\vspace*{\fill}
\thispagestyle{empty}
\newpage
\fi}

%-=-=-=-=-=-=-=-=-=-=-=-=-=-=-=-=-=-=-=-=-=-=-=-=
%	THEOREM STYLES
%-=-=-=-=-=-=-=-=-=-=-=-=-=-=-=-=-=-=-=-=-=-=-=-=

\newcommand{\intoo}[2]{\mathopen{]}#1\,;#2\mathclose{[}}
\newcommand{\ud}{\mathop{\mathrm{{}d}}\mathopen{}}
\newcommand{\intff}[2]{\mathopen{[}#1\,;#2\mathclose{]}}
\newtheorem{notation}{Notation}[chapter]

\newtheoremstyle{ocrenumbox}% % Theorem style name
{0pt}% Space above
{0pt}% Space below
{\normalfont}% % Body font
{}% Indent amount
{\small\bf\sffamily\color{ocre}}% % Theorem head font
{\;}% Punctuation after theorem head
{0.25em}% Space after theorem head
{\small\sffamily\color{ocre}\thmname{#1}\nobreakspace\thmnumber{\@ifnotempty{#1}{}\@upn{#2}}% Theorem text (e.g. Theorem 2.1)
\thmnote{\nobreakspace\the\thm@notefont\sffamily\bfseries\color{black}---\nobreakspace#3.}} % Optional theorem note
\renewcommand{\qedsymbol}{$\blacksquare$}% Optional qed square

\newtheoremstyle{blacknumexa}% Theorem style Example
{5pt}% Space above
{5pt}% Space below
{\normalfont}% Body font
{} % Indent amount
{\small\bf\sffamily\color{sthlmGreen}}% Theorem head font
{\;}% Punctuation after theorem head
{0.25em}% Space after theorem head
{\nobreakspace\thmname{#1}\nobreakspace\thmnumber{\@ifnotempty{#1}{}\@upn{#2}}% Theorem text (e.g. Theorem 2.1)
\thmnote{\nobreakspace\the\thm@notefont\sffamily\bfseries---\nobreakspace#3.}}% Optional theorem note

\newtheoremstyle{blacknumexe}% Theorem style Exercise
{5pt}% Space above
{5pt}% Space below
{\normalfont}% Body font
{} % Indent amount
{\small\bf\sffamily\color{ocre}}% Theorem head font
{\;}% Punctuation after theorem head
{0.25em}% Space after theorem head
{\nobreakspace\thmname{#1}\nobreakspace\thmnumber{\@ifnotempty{#1}{}\@upn{#2}}% Theorem text (e.g. Theorem 2.1)
\thmnote{\nobreakspace\the\thm@notefont\sffamily\bfseries---\nobreakspace#3.}}% Optional theorem note

\newtheoremstyle{blacknumessential}% Theorem style name
{5pt}% Space above
{5pt}% Space below
{\normalfont}% Body font
{} % Indent amount
{\small\bf\sffamily\color{sthlmRed}}% Theorem head font
{\;}% Punctuation after theorem head
{0.25em}% Space after theorem head
{\nobreakspace\thmname{#1}\nobreakspace\thmnumber{\@ifnotempty{#1}{}\@upn{#2}}% Theorem text (e.g. Theorem 2.1)
\thmnote{\nobreakspace\the\thm@notefont\sffamily\bfseries---\nobreakspace#3.}}% Optional theorem note

\newtheoremstyle{blacknumpro}% Theorem style name properties
{5pt}% Space above
{5pt}% Space below
{\normalfont}% Body font
{} % Indent amount
{\small\bf\sffamily\color{sthlmPurple}}% Theorem head font
{\;}% Punctuation after theorem head
{0.25em}% Space after theorem head
{\nobreakspace\thmname{#1}\nobreakspace\thmnumber{\@ifnotempty{#1}{}\@upn{#2}}% Theorem text (e.g. Theorem 2.1)
\thmnote{\nobreakspace\the\thm@notefont\sffamily\bfseries---\nobreakspace#3.}}% Optional theorem note

\newtheoremstyle{blacknumbox} % Theorem style name
{0pt}% Space above
{0pt}% Space below
{\normalfont}% Body font
{}% Indent amount
{\small\bf\sffamily}% Theorem head font
{\;}% Punctuation after theorem head
{0.25em}% Space after theorem head
{\small\sffamily\thmname{#1}\nobreakspace\thmnumber{\@ifnotempty{#1}{}\@upn{#2}}% Theorem text (e.g. Theorem 2.1)
\thmnote{\nobreakspace\the\thm@notefont\sffamily\bfseries---\nobreakspace#3.}}% Optional theorem note

% Non-boxed/non-framed environments
\newtheoremstyle{ocrenum}% % Theorem style name
{5pt}% Space above
{5pt}% Space below
{\normalfont}% % Body font
{}% Indent amount
{\small\bf\sffamily\color{ocre}}% % Theorem head font
{\;}% Punctuation after theorem head
{0.25em}% Space after theorem head
{\small\sffamily\color{ocre}\thmname{#1}\nobreakspace\thmnumber{\@ifnotempty{#1}{}\@upn{#2}}% Theorem text (e.g. Theorem 2.1)
\thmnote{\nobreakspace\the\thm@notefont\sffamily\bfseries\color{black}---\nobreakspace#3.}} % Optional theorem note
\renewcommand{\qedsymbol}{$\blacksquare$}% Optional qed square
\makeatother

% Defines the theorem text style for each type of theorem to one of the three styles above
\newcounter{dummy}
\numberwithin{dummy}{section}
\theoremstyle{ocrenumbox}
\newtheorem{theoremeT}[dummy]{Theorem}
\newtheorem{aruleT}[dummy]{Rule}
\newtheorem{aidentityT}[dummy]{Identity}
\newtheorem{problem}{Problem}[chapter]
\theoremstyle{blacknumexe}
\newtheorem{exerciseT}{Exercise}[chapter]
\theoremstyle{blacknumexa}
\newtheorem{exampleT}{Example}[chapter]
\theoremstyle{blacknumessential}
\newtheorem{essentialqT}{Essential Questions}[chapter]
\theoremstyle{blacknumpro}
\newtheorem{propertyT}[dummy]{Property}
\theoremstyle{blacknumbox}
\newtheorem{notationT}[dummy]{Notation}
\theoremstyle{blacknumbox}
\newtheorem{vocabulary}{Vocabulary}[chapter]
\newtheorem{definitionT}{Definition}[section]
\newtheorem{corollaryT}[dummy]{Corollary}
\theoremstyle{ocrenum}
\newtheorem{proposition}[dummy]{Proposition}

%----------------------------------------------------------------------------------------
%	DEFINITION OF COLORED BOXES
%----------------------------------------------------------------------------------------

\RequirePackage[framemethod=default]{mdframed} % Required for creating the theorem, definition, exercise and corollary boxes

% Theorem box
\newmdenv[skipabove=7pt,
skipbelow=7pt,
backgroundcolor=black!5,
linecolor=ocre,
innerleftmargin=5pt,
innerrightmargin=5pt,
innertopmargin=5pt,
leftmargin=0cm,
rightmargin=0cm,
innerbottommargin=5pt,
linewidth=2pt]{tBox}

% Rule box
\newmdenv[skipabove=7pt,
skipbelow=7pt,
backgroundcolor=black!5,
linecolor=ocre,
innerleftmargin=5pt,
innerrightmargin=5pt,
innertopmargin=5pt,
leftmargin=0cm,
rightmargin=0cm,
innerbottommargin=5pt,
linewidth=2pt]{arBox}

% Identity box
\newmdenv[skipabove=7pt,
skipbelow=7pt,
linecolor=gray,
backgroundcolor=black!5,
innerleftmargin=5pt,
innerrightmargin=5pt,
innertopmargin=5pt,
leftmargin=0cm,
rightmargin=0cm,
innerbottommargin=5pt,
linewidth=2pt]{idenBox}

% Notation box
\newmdenv[skipabove=7pt,
skipbelow=7pt,
backgroundcolor=sthlmOrange!5,
linecolor=sthlmOrange,
innerleftmargin=5pt,
innerrightmargin=5pt,
innertopmargin=5pt,
leftmargin=0cm,
rightmargin=0cm,
innerbottommargin=5pt,
linewidth=2pt]{notBox}

% Definition box
\newmdenv[skipabove=7pt,
skipbelow=7pt,
backgroundcolor=sthlmDarkGrey!5,
linecolor=ocre,
innerleftmargin=5pt,
innerrightmargin=5pt,
innertopmargin=5pt,
innerbottommargin=5pt,
leftmargin=0cm,
rightmargin=0cm,
innerbottommargin=5pt,
linewidth=2pt]{dBox}

% Property box
\newmdenv[skipabove=7pt,
skipbelow=7pt,
backgroundcolor=sthlmDarkGrey!5,
linecolor=sthlmPurple,
innerleftmargin=5pt,
innerrightmargin=5pt,
innertopmargin=5pt,
leftmargin=0cm,
rightmargin=0cm,
innerbottommargin=5pt,
linewidth=2pt]{prBox}

% Exercise box
\newmdenv[skipabove=7pt,
skipbelow=7pt,
rightline=false,
leftline=true,
topline=false,
bottomline=false,
backgroundcolor=ocre!5,
linecolor=ocre,
innerleftmargin=5pt,
innerrightmargin=5pt,
innertopmargin=5pt,
innerbottommargin=5pt,
leftmargin=0cm,
rightmargin=0cm,
linewidth=4pt]{eBox}

% Example box
\newmdenv[skipabove=7pt,
skipbelow=7pt,
rightline=false,
leftline=true,
topline=false,
bottomline=false,
backgroundcolor=sthlmGreen!5,
linecolor=sthlmGreen,
innerleftmargin=5pt,
innerrightmargin=5pt,
innertopmargin=5pt,
innerbottommargin=5pt,
leftmargin=0cm,
rightmargin=0cm,
linewidth=4pt]{exaBox}

% Essential Question box
\newmdenv[skipabove=7pt,
skipbelow=7pt,
rightline=false,
leftline=false,
topline=true,
bottomline=true,
backgroundcolor=sthlmRed!5,
linecolor=sthlmRed,
innerleftmargin=5pt,
innerrightmargin=5pt,
innertopmargin=5pt,
innerbottommargin=5pt,
leftmargin=0cm,
rightmargin=0cm,
linewidth=4pt]{essentialqBox}

% Corollary box
\newmdenv[skipabove=7pt,
skipbelow=7pt,
rightline=false,
leftline=true,
topline=false,
bottomline=false,
linecolor=gray,
backgroundcolor=black!5,
innerleftmargin=5pt,
innerrightmargin=5pt,
innertopmargin=5pt,
leftmargin=0cm,
rightmargin=0cm,
linewidth=4pt,
innerbottommargin=5pt]{cBox}

% Creates an environment for each type of theorem and assigns it a theorem text style from the "Theorem Styles" section above and a colored box from above
\newenvironment{theorem}{\begin{tBox}\begin{theoremeT}}{\end{theoremeT}\end{tBox}}
\newenvironment{arule}{\begin{arBox}\begin{aruleT}}{\end{aruleT}\end{arBox}}
\newenvironment{aidentity}{\begin{idenBox}\begin{aidentityT}}{\end{aidentityT}\end{idenBox}}
\newenvironment{exercise}{\begin{eBox}\begin{exerciseT}}{\hfill{\color{ocre}\tiny\ensuremath{\blacksquare}}\end{exerciseT}\end{eBox}}
\newenvironment{definition}{\begin{dBox}\begin{definitionT}}{\end{definitionT}\end{dBox}}
\newenvironment{property}{\begin{prBox}\begin{propertyT}}{\end{propertyT}\end{prBox}}
\newenvironment{notations}{\begin{notBox}\begin{notationT}}{\end{notationT}\end{notBox}}
\newenvironment{example}{\begin{exaBox}\begin{exampleT}}{\hfill{\color{sthlmGreen}\tiny\ensuremath{\blacksquare}}\end{exampleT}\end{exaBox}}
\newenvironment{essentialq}{\begin{essentialqBox}\begin{essentialqT}}{\end{essentialqT}\end{essentialqBox}}
\newenvironment{corollary}{\begin{cBox}\begin{corollaryT}}{\end{corollaryT}\end{cBox}}

%----------------------------------------------------------------------------------------
%	REMARK ENVIRONMENT
%----------------------------------------------------------------------------------------

\newenvironment{remark}{\par\vspace{10pt}\small % Vertical white space above the remark and smaller font size
\begin{list}{}{
\leftmargin=35pt % Indentation on the left
\rightmargin=25pt}\item\ignorespaces % Indentation on the right
\makebox[-2.5pt]{\begin{tikzpicture}[overlay]
\node[draw=ocre!60,line width=1pt,circle,fill=ocre!25,font=\sffamily\bfseries,inner sep=2pt,outer sep=0pt] at (-15pt,0pt){\textcolor{ocre}{R}};\end{tikzpicture}} % Orange R in a circle
\advance\baselineskip -1pt}{\end{list}\vskip5pt} % Tighter line spacing and white space after remark

\newcommand{\soln}[1]{\par\vspace{10pt}\makebox[15pt]{\begin{tikzpicture}[overlay]
\node[draw=sthlmGreen!70,line width=1pt,circle,fill=sthlmGreen!25,font=\sffamily\bfseries,inner sep=2pt,outer sep=0pt] at (-15pt,0pt){\textcolor{sthlmGreen}{S}};\end{tikzpicture}}\cGreen{\noindent\hrulefill\vskip 15pt}}

\newcommand{\qdepend}[1]{\par\vspace{10pt}\makebox[15pt]{\begin{tikzpicture}[overlay]
\node[draw=sthlmDarkGrey!70,line width=1pt,circle,fill=sthlmDarkGrey!25,font=\sffamily\bfseries,inner sep=2pt,outer sep=0pt] at (-15pt,0pt){\textcolor{sthlmDarkGrey}{D}};\end{tikzpicture}}\cDarkGrey{\noindent\hrulefill\vskip 15pt}}

%-=-=-=-=-=-=-=-=-=-=-=-=-=-=-=-=-=-=-=-=-=-=-=-=
%	SECTION NUMBERING IN THE MARGIN
%-=-=-=-=-=-=-=-=-=-=-=-=-=-=-=-=-=-=-=-=-=-=-=-=

\makeatletter
\renewcommand{\@seccntformat}[1]{\llap{\textcolor{ocre}{\csname the#1\endcsname}\hspace{1em}}}
\renewcommand{\section}{\@startsection{section}{1}{\z@}
{-4ex \@plus -1ex \@minus -.4ex}
{1ex \@plus.2ex }
{\normalfont\large\sffamily\bfseries}}
\renewcommand{\subsection}{\@startsection {subsection}{2}{\z@}
{-3ex \@plus -0.1ex \@minus -.4ex}
{0.5ex \@plus.2ex }
{\normalfont\sffamily\bfseries}}
\renewcommand{\subsubsection}{\@startsection {subsubsection}{3}{\z@}
{-2ex \@plus -0.1ex \@minus -.2ex}
{.2ex \@plus.2ex }
{\normalfont\small\sffamily\bfseries}}
\renewcommand\paragraph{\@startsection{paragraph}{4}{\z@}
{-2ex \@plus-.2ex \@minus .2ex}
{.1ex}
{\normalfont\small\sffamily\bfseries}}

%-=-=-=-=-=-=-=-=-=-=-=-=-=-=-=-=-=-=-=-=-=-=-=-=
%	PART HEADINGS
%-=-=-=-=-=-=-=-=-=-=-=-=-=-=-=-=-=-=-=-=-=-=-=-=

% numbered part in the table of contents
\newcommand{\@mypartnumtocformat}[2]{%
\setlength\fboxsep{0pt}%
\noindent\colorbox{ocre!20}{\strut\parbox[c][.7cm]{\ecart}{\color{ocre!70}\Large\sffamily\bfseries\centering#1}}\hskip\esp\colorbox{ocre!40}{\strut\parbox[c][.7cm]{\linewidth-\ecart-\esp}{\Large\sffamily\centering#2}}}%
%%%%%%%%%%%%%%%%%%%%%%%%%%%%%%%%%%
% unnumbered part in the table of contents
\newcommand{\@myparttocformat}[1]{%
\setlength\fboxsep{0pt}%
\noindent\colorbox{ocre!40}{\strut\parbox[c][.7cm]{\linewidth}{\Large\sffamily\centering#1}}}%
%%%%%%%%%%%%%%%%%%%%%%%%%%%%%%%%%%
\newlength\esp
\setlength\esp{4pt}
\newlength\ecart
\setlength\ecart{1.2cm-\esp}
\newcommand{\thepartimage}{}%
\newcommand{\partimage}[1]{\renewcommand{\thepartimage}{#1}}%
\def\@part[#1]#2{%
\ifnum \c@secnumdepth >-2\relax%
\refstepcounter{part}%
\addcontentsline{toc}{part}{\texorpdfstring{\protect\@mypartnumtocformat{\thepart}{#1}}{\partname~\thepart\ ---\ #1}}
\else%
\addcontentsline{toc}{part}{\texorpdfstring{\protect\@myparttocformat{#1}}{#1}}%
\fi%
\startcontents%
\markboth{}{}%
{\thispagestyle{empty}%
\begin{tikzpicture}[remember picture,overlay]%
\node at (current page.north west){\begin{tikzpicture}[remember picture,overlay]%
\fill[ocre!20](0cm,0cm) rectangle (\paperwidth,-\paperheight);
\node[anchor=north] at (4cm,-3.25cm){\color{ocre!40}\fontsize{220}{100}\sffamily\bfseries\@Roman\c@part};
\node[anchor=south east] at (\paperwidth-1cm,-\paperheight+1cm){\parbox[t][][t]{8.5cm}{
\printcontents{l}{0}{\setcounter{tocdepth}{1}}%
}};
\node[anchor=north east] at (\paperwidth-1.5cm,-3.25cm){\parbox[t][][t]{15cm}{\strut\raggedleft\color{white}\fontsize{30}{30}\sffamily\bfseries#2}};
\end{tikzpicture}};
\end{tikzpicture}}%
\@endpart}
\def\@spart#1{%
\startcontents%
\phantomsection
{\thispagestyle{empty}%
\begin{tikzpicture}[remember picture,overlay]%
\node at (current page.north west){\begin{tikzpicture}[remember picture,overlay]%
\fill[ocre!20](0cm,0cm) rectangle (\paperwidth,-\paperheight);
\node[anchor=north east] at (\paperwidth-1.5cm,-3.25cm){\parbox[t][][t]{15cm}{\strut\raggedleft\color{white}\fontsize{30}{30}\sffamily\bfseries#1}};
\end{tikzpicture}};
\end{tikzpicture}}
\addcontentsline{toc}{part}{\texorpdfstring{%
\setlength\fboxsep{0pt}%
\noindent\protect\colorbox{ocre!40}{\strut\protect\parbox[c][.7cm]{\linewidth}{\Large\sffamily\protect\centering #1\quad\mbox{}}}}{#1}}%
\@endpart}
\def\@endpart{\vfil\newpage
\if@twoside
\if@openright
\null
\thispagestyle{empty}%
\newpage
\fi
\fi
\if@tempswa
\twocolumn
\fi}

%-=-=-=-=-=-=-=-=-=-=-=-=-=-=-=-=-=-=-=-=-=-=-=-=
%	CHAPTER HEADINGS
%-=-=-=-=-=-=-=-=-=-=-=-=-=-=-=-=-=-=-=-=-=-=-=-=

%\newcommand{\thechapterimage}{}%
%\newcommand{\chapterimage}[1]{\renewcommand{\thechapterimage}{#1}}%
\def\@makechapterhead#1{%
{\parindent \z@ \raggedright \normalfont
\ifnum \c@secnumdepth >\m@ne
\if@mainmatter
\begin{tikzpicture}[remember picture,overlay]
\node at (current page.north west)
{\begin{tikzpicture}[remember picture,overlay]
%\node[anchor=north west,inner sep=0pt] at (0,0) {\includegraphics[width=\paperwidth]{\thechapterimage}};
\draw[anchor=west] (\Gm@lmargin,-4cm) node [line width=2pt,rounded corners=15pt,draw=ocre,fill=white,fill opacity=0.5,inner sep=15pt]{\strut\makebox[22cm]{}};
\draw[anchor=west] (\Gm@lmargin+.3cm,-4cm) node {\huge\sffamily\bfseries\color{black}\thechapter. #1\strut};
\end{tikzpicture}};
\end{tikzpicture}
\else
\begin{tikzpicture}[remember picture,overlay]
\node at (current page.north west)
{\begin{tikzpicture}[remember picture,overlay]
%\node[anchor=north west,inner sep=0pt] at (0,0) {\includegraphics[width=\paperwidth]{\thechapterimage}};
\draw[anchor=west] (\Gm@lmargin,-4cm) node [line width=2pt,rounded corners=15pt,draw=ocre,fill=white,fill opacity=0.5,inner sep=15pt]{\strut\makebox[22cm]{}};
\draw[anchor=west] (\Gm@lmargin+.3cm,-4cm) node {\huge\sffamily\bfseries\color{black}#1\strut};
\end{tikzpicture}};
\end{tikzpicture}
\fi\fi\par\vspace*{100\p@}}}

%-------------------------------------------

\def\@makeschapterhead#1{%
\begin{tikzpicture}[remember picture,overlay]
\node at (current page.north west)
{\begin{tikzpicture}[remember picture,overlay]
%\node[anchor=north west,inner sep=0pt] at (0,0) {\includegraphics[width=\paperwidth]{\thechapterimage}};
\draw[anchor=west] (\Gm@lmargin,-4cm) node [line width=2pt,rounded corners=15pt,draw=ocre,fill=white,fill opacity=0.5,inner sep=15pt]{\strut\makebox[22cm]{}};
\draw[anchor=west] (\Gm@lmargin+.3cm,-4cm) node {\huge\sffamily\bfseries\color{black}#1\strut};
\end{tikzpicture}};
\end{tikzpicture}
\par\vspace*{100\p@}}
\makeatother

%-=-=-=-=-=-=-=-=-=-=-=-=-=-=-=-=-=-=-=-=-=-=-=-=
%	HYPERLINKS IN THE DOCUMENTS
%-=-=-=-=-=-=-=-=-=-=-=-=-=-=-=-=-=-=-=-=-=-=-=-=

\usepackage{hyperref}
\hypersetup{hidelinks,backref=true,pagebackref=true,hyperindex=true,colorlinks=false,breaklinks=true,urlcolor= ocre,bookmarks=true,bookmarksopen=false,pdftitle={Title},pdfauthor={Author}}
\usepackage{bookmark}
\bookmarksetup{
open,
numbered,
addtohook={%
\ifnum\bookmarkget{level}=0 % chapter
\bookmarksetup{bold}%
\fi
\ifnum\bookmarkget{level}=-1 % part
\bookmarksetup{color=ocre,bold}%
\fi
}
}

%-=-=-=-=-=-=-=-=-=-=-=-=-=-=-=-=-=-=-=-=-=-=-=-=
%	TIKZ LIBRARIES
%-=-=-=-=-=-=-=-=-=-=-=-=-=-=-=-=-=-=-=-=-=-=-=-=

\usetikzlibrary{
arrows,
backgrounds,
fit,
chains,
calc,
decorations.pathmorphing,
matrix,
mindmap,
positioning,
shapes,
decorations,
shadows,
trees
}

\tikzset{ >=stealth', help lines/.style={dashed, thick}, axis/.style={<->}, important line/.style={thick}, connection/.style={thick, dotted},}

%-=-=-=-=-=-=-=-=-=-=-=-=-=-=-=-=-=-=-=-=-=-=-=-=
%	TIKZ FACTOR
%-=-=-=-=-=-=-=-=-=-=-=-=-=-=-=-=-=-=-=-=-=-=-=-=
\tikzstyle{firstterm} = [circle, draw, fill=sthlmRed!20, text centered,minimum size=1cm]
\tikzstyle{secondterm} = [circle, draw, fill=sthlmBlue!20, text centered, node distance=0.25cm,minimum size=1cm]
\tikzstyle{terms} = [rectangle, draw, fill=sthlmDarkGrey!20, text centered, node distance=0.25cm,minimum size=1cm]
\tikzstyle{factoradd} = [rectangle, draw=none, text centered, node distance=0.25cm,minimum size=1cm]
\tikzstyle{multiply} = [rectangle, draw, fill=sthlmDarkGrey!20, text centered, minimum height=0.75cm,minimum size=1cm]
\tikzstyle{add} = [rectangle, draw, fill=sthlmGreen!20, text centered, minimum height=0.75cm,minimum size=1cm]
\tikzstyle{line} = [draw, color=black!50, -latex']
% Define the style for the red dotted boxes
\tikzset{factordotted/.style={draw=sthlmPurple!50!white, line width=1pt, dash pattern=on 1pt off 4pt on 6pt off 4pt, inner sep=4mm, rectangle, rounded corners}}

%-=-=-=-=-=-=-=-=-=-=-=-=-=-=-=-=-=-=-=-=-=-=-=-=
%	WORKFLOW SETUP: SIMPLIFYING EXPRESSIONS
%-=-=-=-=-=-=-=-=-=-=-=-=-=-=-=-=-=-=-=-=-=-=-=-=

\tikzstyle{operation} = [rectangle, draw, fill=sthlmBlue!40, text centered, rounded corners]
\tikzstyle{function} = [rectangle, draw, fill=sthlmPurple!40, text centered, rounded corners]
\tikzstyle{notation} = [rectangle, draw, fill=sthlmGreen!40, text centered, rounded corners]
\tikzstyle{algorithm} = [rectangle, draw, fill=sthlmOrange!40, text centered, rounded corners]
\tikzstyle{delim} = [rectangle, draw, fill=sthlmYellow!40, text centered, rounded corners]
\tikzstyle{property} = [rectangle, draw, fill=sthlmRed!40, text centered, rounded corners]
\tikzstyle{hide} = [rectangle, draw, fill=sthlmGrey, text centered, rounded corners]
\tikzstyle{active} = [rectangle, draw, fill=sthlmDarkGrey, text=white, text centered, rounded corners]

\tikzset{functiondotted/.style={draw=sthlmPurple!50!white, line width=1pt,
                               dash pattern=on 1pt off 4pt on 6pt off 4pt,
                                inner sep=4mm, rectangle, rounded corners}}

\tikzset{productdotted/.style={draw=sthlmBlue!50!white, line width=1pt,
                               dash pattern=on 1pt off 4pt on 6pt off 4pt,
                                inner sep=4mm, rectangle, rounded corners}}

\tikzset{sumdotted/.style={draw=sthlmRed!50!white, line width=1pt,
                               dash pattern=on 1pt off 4pt on 6pt off 4pt,
                                inner sep=4mm, rectangle, rounded corners}}

%-=-=-=-=-=-=-=-=-=-=-=-=-=-=-=-=-=-=-=-=-=-=-=-=
%	WORKFLOW SETUP TRIGONOMETRIC EQUATIONS
%-=-=-=-=-=-=-=-=-=-=-=-=-=-=-=-=-=-=-=-=-=-=-=-=

\tikzstyle{decision} = [diamond, draw, fill=sthlmGreen!20,
    text width=4.5em, text badly centered, node distance=1.2cm, inner sep=0pt]
\tikzstyle{abscissa} = [rectangle, draw, fill=sthlmBlue!20,
    text width=8em, text centered, rounded corners, node distance=1.2cm, minimum height=2.8em]
\tikzstyle{ordinate} = [rectangle, draw, fill=sthlmRed!20,
    text width=8em, text centered, rounded corners, node distance=1.2cm, minimum height=2.8em]
\tikzstyle{task} = [rectangle, draw, fill=sthlmPurple!20,
   minimum width=8em, text centered, rounded corners, node distance=1.2cm, minimum height=2.8em]
\tikzstyle{pointcircle} = [circle, draw, fill=sthlmDarkGrey!20, text centered, node distance=1cm]
% Define the style for the red dotted boxes
  \tikzset{bluedotted/.style={draw=sthlmBlue!50!white, line width=1pt,
                               dash pattern=on 1pt off 4pt on 6pt off 4pt,
                                inner sep=4mm, rectangle, rounded corners}}
% Define the style for the red dotted boxes
  \tikzset{reddotted/.style={draw=sthlmRed!50!white, line width=1pt,
                               dash pattern=on 1pt off 4pt on 6pt off 4pt,
                                inner sep=4mm, rectangle, rounded corners}}

%-=-=-=-=-=-=-=-=-=-=-=-=-=-=-=-=-=-=-=-=-=-=-=-=
%	Derivative Workflow
%-=-=-=-=-=-=-=-=-=-=-=-=-=-=-=-=-=-=-=-=-=-=-=-=	
\tikzstyle{point} = [rectangle, draw, fill=sthlmBlue!40, text centered,rounded corners]
\tikzstyle{limit} = [rectangle, draw, fill=sthlmPurple!40, text centered,rounded corners]
\tikzstyle{class} = [rectangle, draw, fill=sthlmGreen!40, text centered, rounded corners]
\tikzstyle{intervals} = [rectangle, draw, fill=sthlmOrange!40, text centered, rounded corners]
\tikzstyle{fd} = [rectangle, draw, fill=sthlmRed!40, text centered, rounded corners]
\tikzstyle{secondderivative} = [circle, draw, fill=sthlmGrey,radius=0.5, text width=1em, node distance =2.5cm]

%-=-=-=-=-=-=-=-=-=-=-=-=-=-=-=-=-=-=-=-=-=-=-=-=
%	General Workflow
%-=-=-=-=-=-=-=-=-=-=-=-=-=-=-=-=-=-=-=-=-=-=-=-=	
\tikzstyle{bluerbox} = [rectangle, draw, fill=sthlmBlue!40, text centered,rounded corners]
\tikzstyle{purplerbox} = [rectangle, draw, fill=sthlmPurple!40, text centered,rounded corners]
\tikzstyle{greenrbox} = [rectangle, draw, fill=sthlmGreen!40, text centered, rounded corners]
\tikzstyle{orangerbox} = [rectangle, draw, fill=sthlmOrange!40, text centered, rounded corners]
\tikzstyle{redrbox} = [rectangle, draw, fill=sthlmRed!40, text centered, rounded corners]
\tikzstyle{circlenode} = [circle, draw, fill=sthlmGrey,radius=0.5, text width=1em, node distance =2.5cm]
\tikzstyle{line} = [draw, very thick, color=black!50, -latex']
\tikzstyle{dashline} = [draw, very thick, dash pattern=on 1pt off 4pt on 6pt off 4pt,inner sep=4mm,color=black!50, -latex']

\begin{document}
%-=-=-=-=-=-=-=-=-=-=-=-=-=-=-=-=-=-=-=-=-=-=-=-=
%	TITLE PAGE
%-=-=-=-=-=-=-=-=-=-=-=-=-=-=-=-=-=-=-=-=-=-=-=-=

\begingroup
\thispagestyle{empty}
\begin{tikzpicture}[remember picture,overlay]
\coordinate [below=12cm] (midpoint) at (current page.north);
\node at (current page.north west)
{\begin{tikzpicture}[remember picture,overlay]
\node[anchor=north west,inner sep=0pt] at (0,0) {\includegraphics[width=\paperwidth]{Pictures/background}}; % Background image
\draw[anchor=north] (midpoint) node [fill=ocre!30!white,fill opacity=0.6,text opacity=1,inner sep=1cm]{\Huge\centering\bfseries\sffamily\parbox[c][][t]{\paperwidth}{\centering \cDarkGrey{Mathematics Notebook}\\[15pt] % Book title
{\cDarkGrey{\huge by Mark Olson : \texttt{v42.org}}}\\[20pt] % Subtitle
{\cDarkGrey{\large Updated: \today}}}}; 
\end{tikzpicture}};
\end{tikzpicture}
\vfill
\endgroup

%-=-=-=-=-=-=-=-=-=-=-=-=-=-=-=-=-=-=-=-=-=-=-=-=
%	COPYRIGHT PAGE
%-=-=-=-=-=-=-=-=-=-=-=-=-=-=-=-=-=-=-=-=-=-=-=-=

\newpage
%~\vfill
%\thispagestyle{empty}

%\noindent Copyright \copyright\ 2013 John Smith\\ % Copyright notice

%\noindent \textsc{Published by Publisher}\\ % Publisher

%\noindent \textsc{book-website.com}\\ % URL

%\noindent Licensed under the Creative Commons Attribution-NonCommercial 3.0 Unported License (the ``License''). You may not use this file except in compliance with the License. You may obtain a copy of the License at \url{http://creativecommons.org/licenses/by-nc/3.0}. Unless required by applicable law or agreed to in writing, software distributed under the License is distributed on an \textsc{``as is'' basis, without warranties or conditions of any kind}, either express or implied. See the License for the specific language governing permissions and limitations under the License.\\ % License information

%\noindent \textit{First printing, March 2013} % Printing/edition date

%-=-=-=-=-=-=-=-=-=-=-=-=-=-=-=-=-=-=-=-=-=-=-=-=
%	TABLE OF CONTENTS
%-=-=-=-=-=-=-=-=-=-=-=-=-=-=-=-=-=-=-=-=-=-=-=-=

%\chapterimage{Pictures/chapter_head_1.pdf} % Table of contents heading image

\pagestyle{empty} % No headers

\tableofcontents % Print the table of contents itself

\cleardoublepage % Forces the first chapter to start on an odd page so it's on the right

\pagestyle{fancy} % Print headers again

%-=-=-=-=-=-=-=-=-=-=-=-=-=-=-=-=-=-=-=-=-=-=-=-=
%	PARTS
%-=-=-=-=-=-=-=-=-=-=-=-=-=-=-=-=-=-=-=-=-=-=-=-=
\part{Reference}
\subfile{c-numbers}
\subfile{c-operations}
\subfile{c-notation}
\subfile{c-properties}
\subfile{c-identities}
\subfile{c-simplifyworkflow}
\subfile{c-calculusrules}

\part{Algebra}
\subfile{algebra-simplify-univariatepolynomials}
\subfile{algebra-simplify-multivariatemonomials}
\subfile{algebra-solve-deg1polynomials}
\subfile{algebra-solve-deg2polynomials}

\part{Functions}
\subfile{functions-evaluate-polynomials}
\part{Differential Calculus}
\subfile{calculus-differentiation-firstprinciples}
\subfile{calculus-differentiation-rules}
\subfile{calculus-differentiation-equationtangentsecantlines}
\subfile{calculus-differentiation-1stderivativetest}
\subfile{calculus-differentiation-localextrema}
\subfile{calculus-differentiation-2ndderivativetest}
\subfile{calculus-differentiation-curvesketching}
%\part{MYP}
%\subfile{chapter-myp-algebraicexpressions}
%\subfile{chapter-myp-arithmeticexpressions}
%\documentclass[20150903-160354-rs2.2-MarksMathNotebook.tex]{subfiles}

\begin{document}
%-=-=-=-=-=-=-=-=-=-=-=-=-=-=-=-=-=-=-=-=-=-=-=-=
%
%	CHAPTER
%
%-=-=-=-=-=-=-=-=-=-=-=-=-=-=-=-=-=-=-=-=-=-=-=-=

\chapterimage{Pictures/chapter_head_2.pdf} % Chapter heading image

\chapter{Algebraic Expressions}

%-=-=-=-=-=-=-=-=-=-=-=-=-=-=-=-=-=-=-=-=-=-=-=-=
%	SECTION: OPERATION OF ADDITION
%-=-=-=-=-=-=-=-=-=-=-=-=-=-=-=-=-=-=-=-=-=-=-=-=
\section{Expressions}\index{Algebraic Expressions}

\begin{essentialq}\hfill \\

\begin{enumerate}
	\item What is an algebraic expression?
\end{enumerate}

\end{essentialq}


\begin{table}
\begin{tabular}{|l|c|c|c|c|c|c}
\hline
 									& Arithmetic 	& Polynomial	& Algebraic  	\\
\hline
Constant 							& \cellyes		& \cellyes 		& \cellyes		\\
\hline
Factorial 							& \cellyes 		& \cellyes 		& \cellyes		\\
\hline
Variable: parameter/coefficient 	& \cellyes 		& \cellyes 		& \cellyes		\\
\hline
Variable: unknown/indeterminate 	& \cellno 		& \cellyes 		& \cellyes		\\
\hline
Power with $\mathbb{Z}^{+}$ exponent 	& \cellno 		& \cellyes 		& \cellyes		 \\
\hline
Power with $\mathbb{Z}$ exponent 		& \cellno 		& \cellno 		& \cellyes		 \\
\hline
$n$-th root 						& \cellno 		& \cellno 		& \cellyes		\\
\hline
Power with $\mathbb{Q}$ exponent 	& \cellno 		& \cellno 		& \cellyes		 \\
\hline
\end{tabular}
\caption{Names of different types of expressions}\label{tab:tableofexpressions}
\end{table}


%-=-=-= DEFINITION
%\begin{definition}[Algebraic Expression]\index{Algebraic Expression}

%An algebraic expression is made up of one or more terms (operation of addition), which are themselves made up of factors (operation of multiplication).\\

%Algebraic expressions have three different types of multiplicative factors:\\

%\begin{enumerate}
%	\item constant (coefficient)
%	\begin{itemize}
%		\item real number
%		\item parameter
%	\end{itemize}
%	\item variable
%	\item rational power
%\end{enumerate}
%\hfill \cite{mathworld:algebraicexpression}
%\end{definition}

\section{Polynomial Expressions}\index{Polynomial Expressions}






\newpage
\begin{figure}
\begin{tikzpicture}[scale=1, node distance=1.1cm, text width=5em, auto]
    % Place nodes
\node [property] (MId1) {MId};
\node [notation, below of=MId1] (ONeg1) {ONeg};
\node [operation, below of=ONeg1] (DOS1) {DOS};
\node [delim, below of=DOS1] (DELIM) {DELIM};
\node [property, below of=DELIM] (DPE) {DPE};
\node [notation, below of=DPE] (JTC) {JTC};
\node [property, below of=JTC] (CPM) {CPM};
\node [operation, below of=CPM] (OOM) {OOM};
\node [property, right of=OOM, node distance=2.8cm] (APM) {APM};
\node [algorithm, below of=OOM] (RF1) {RF};
\node [notation, below of=RF1] (CTJ) {CTJ};
\node [notation, below of=CTJ] (CPA) {CPA};
\node [property, below of=CPA] (DPF) {DPF};
\node [operation, below of=DPF] (OOA) {OOA};
\node [property, right of=OOA,node distance=2.8cm] (APA) {APA};
\node [algorithm, below of=OOA] (RF2) {RF};
\node [property, below of=RF2] (AId) {AId};
\node [operation, below of=AId] (DOS2) {DOS};
\node [notation, below of=DOS2] (ONeg2) {ONeg};
\node [property, below of=ONeg2] (MId2) {MId};

\path [line] (MId1) -- (ONeg1);
\path [line] (ONeg1) -- (DOS1);
\path [line] (DOS1) -- (DELIM);
\path [line] (DELIM) -- (DPE);
\path [line] (DPE) -- (JTC);
\path [line] (JTC) -- (CPM);
\path [line] (CPM) edge [out= 315, in= 150] (APM);
\path [line] (APM) -- (OOM);
\path [line] (OOM) -- (RF1);
\path [line] (RF1) -- (CTJ);
\path [line] (CTJ) -- (CPA);
\path [line] (CPA) -- (DPF);
\path [line] (DPF) edge [out= 315, in= 150] (APA);
\path [line] (APA) -- (OOA);
\path [line] (OOA) -- (RF2);
\path [line] (RF2) -- (AId);
\path [line] (AId) -- (DOS2);
\path [line] (DOS2) -- (ONeg2);
\path [line] (ONeg2) -- (MId2);

\node [function, right of=DELIM, node distance=8cm] (PoTF) {PoTF};
\node [function, below of=PoTF] (RTPo) {RTPo};
\node [function, below of=RTPo] (PoNE1) {PoNE};
\node [function, below of=PoNE1] (PoPr) {PoPr};
\node [function, below of=PoPr] (PoQ) {PoQ};
\node [function, below of=PoQ] (PoPrPo) {PoPrPo};
\node [function, below of=PoPrPo] (PoQPo) {PoQPo};
\node [function, below of=PoQPo] (PrCBPo) {PrCBPo};
\node [function, below of=PrCBPo] (QCBPo) {QCBPo};
\node [function, below of=QCBPo] (PoPo) {PoPo};
\node [function, below of=PoPo] (PoNE2) {PoNE};
\node [operation, below of=PoNE2] (OOE) {OOE};
\node [function, below of=OOE] (PoTR) {PoTR};


\path [line] (DELIM) edge [bend left=30] (PoTF);
\path [line] (PoTF) -- (RTPo);
\path [line] (RTPo) -- (PoNE1);
\path [line] (PoPr) -- (PoQ);
\path [line] (PoQ) -- (PoPrPo);
\path [line] (PoPrPo) -- (PoQPo);
\path [line] (PoQPo) -- (PrCBPo);
\path [line] (PrCBPo) -- (QCBPo);
\path [line] (QCBPo) -- (PoPo);
\path [line] (PoPo) -- (PoNE2);
\path [line] (PoNE2) -- (OOE);
\path [line] (OOE) -- (PoTR);

%\node[productdotted, fit =(DELIM)(JTC)(CPM)(OOM)(RF1)(CTJ)(APM)](product) {};
%\node[left of =product, text centered,](powers) {\textcolor{sthlmPurple}{\textbf{Multiplication}}};
%\node[sumdotted, fit =(CPA)(DPF)(OOA)(RF2)(AId)(ONeg2)(DOS2)(MId2)(APA)](sum2) {};
%\node[sumdotted, fit =(MId1)(ONeg1)(DOS1)](sum1) {};
%\node[left of =sum1, text centered,](powers) {\textcolor{sthlmPurple}{\textbf{Addition}}};
%\node[left of =sum2, text centered,](powers) {\textcolor{sthlmPurple}{\textbf{Addition}}};


\node[functiondotted, fit =(PoTF)(RTPo)(PoNE1)(PoPr)(PoQ)(PoPrPo)(PoQPo)(PrCBPo)(QCBPo)(PoPo)(PoNE2)(OOE)(PoTR)](powerfunctions) {};
\node[below of =PoTR, text centered,](powers) {\textcolor{sthlmPurple}{\textbf{Powers}}};

\path [line] (PoNE1) edge [out= 180, in= 0] (DELIM);

\path [line] (CPM) edge [out= 0, in= 180] (PoPr);

\path [line] (PoTR) edge [out= 180, in= 0] (APM);

\path [line] (AId) edge [out= 180, in= 180] (DELIM);

\end{tikzpicture}
\caption{Simplifying Expressions Workflow: \newline {\color{sthlmRed!40}\ensuremath{\blacksquare}} Property, {\color{sthlmBlue!40}\ensuremath{\blacksquare}} Operation, {\color{sthlmGreen!40}\ensuremath{\blacksquare}} Notation, {\color{sthlmPurple!40}\ensuremath{\blacksquare}} Powers, {\color{sthlmYellow!40}\ensuremath{\blacksquare}} Delimiters, {\color{sthlmOrange!40}\ensuremath{\blacksquare}} Process, {\color{sthlmDarkGrey!40}\ensuremath{\blacksquare}} Not Used}
\end{figure}






%-=-=-= DEFINITION
\begin{definition}[Operation of Exponentiation (OOE)]\index{Operation!Operation of Exponentiation}
\begin{align}
\underbrace{{\underbrace{b}_{base}}^{\overbrace{m}^{Exponent}}}_{Power} \label{eq:ooe}
\end{align}
\end{definition}


%-=-=-=-=-=-=-=-=-=-=-=-=-=-=-=-=-=-=-=-=-=-=-=-=
%	SECTION: ALGEBRAIC EXPRESSIONS
%-=-=-=-=-=-=-=-=-=-=-=-=-=-=-=-=-=-=-=-=-=-=-=-=






%-=-=-= DEFINITION
\begin{definition}[Juxtaposition to Center-Dot (JTC)]\index{Notation!Juxtapostion to Center-Dot}
\begin{align}
ab &= a \cdot b \label{eq:jtc}
\end{align}
\end{definition}

%-=-=-= DEFINITION
\begin{definition}[Center-Dot to Justapostion (CTJ)]\index{Notation!Center-Dot to Juxtaposition}
\begin{align}
a \cdot b &= ab \label{eq:ctj}
\end{align}
\end{definition}


%-=-=-= DEFINITION
\begin{definition}[Commutative Property of Multiplication (CPM)]\index{Property!Commutative Property of Multiplication}
\begin{align}
\alert{a} \cdot b &= b \cdot \alert{a} \label{eq:cpm}
\end{align}
\end{definition}



%-=-=-= DEFINITION
\begin{definition}[Multiplicative Inverse (MI)]\index{Multiplicative Inverse}
\begin{subequations}
\begin{align}
a \cdot \alert{\dfrac{1}{a}} &= 1 \label{eq:mi1} \\
a \cdot \alert{a^{-1}} &= 1 \label{eq:mi2}
\end{align}
\end{subequations}
\end{definition}

%-=-=-= DEFINITION
\begin{definition}[Associative Property of Multiplication (APM)]\index{Property!Associative Property of Multiplication}
\begin{subequations}
\begin{align}
a\cdot b\cdot c &= (a\cdot b)\cdot c \label{eq:apm1} \\
a\cdot b\cdot c &= a\cdot (b\cdot c) \label{eq:apm2}
\end{align}
\end{subequations}
\end{definition}


%-=-=-=-=-=-=-=-=-=-=-=-=-=-=-=-=-=-=-=-=-=-=-=-=
%	SECTION: Powers
%-=-=-=-=-=-=-=-=-=-=-=-=-=-=-=-=-=-=-=-=-=-=-=-=

\subsubsection{Powers}\index{Powers}


\begin{arule}[Power of a Quotient of Powers (PoQPo)]\index{Powers!Power of a Quotient of Powers}
\begin{subequations}
\begin{align}
	\left(\dfrac{a^m}{b^n}\right)^k &= \dfrac{a^{m \cdot k}}{b^{n\cdot k}} \label{eq:poqpo1}\\
	\dfrac{a^{m \cdot k}}{b^{n\cdot k}} &= \left(\dfrac{a^m}{b^n}\right)^k \label{eq:poqpo2}
\end{align}
\end{subequations}
\end{arule}

\begin{arule}[Power of a Product of Powers (PoPrPo)]\index{Powers!Power of a Product of Powers}
\begin{subequations}
\begin{align}
	\left(a^m \cdot b^n\right)^k &= a^{m \cdot k} \cdot b^{n\cdot k} \label{eq:poprpo1}\\
	a^{m \cdot k} \cdot b^{n \cdot k} &= \left(a^m \cdot b^n\right)^k \label{eq:poprpo2}
\end{align}
\end{subequations}
\end{arule}

\begin{definition}[Power To Factor (PoTF)]\index{Powers!Power to Factor}
\begin{align}
a^{\alert{n}} &= a_1 \cdot a_2 \cdot \ldots \cdot a_{n-1} \cdot a_{\alert{n}} \label{eq:potf}
\end{align}
\end{definition}

\begin{definition}[Factor To Power (FTPo)]\index{Powers!Factor to Power}
\begin{align}
a_1 \cdot a_2 \cdot \ldots \cdot a_{n-1} \cdot a_{\alert{n}} &= a^{\alert{n}} \label{eq:ftpo}
\end{align}
\end{definition}

%-=-=-= DEFINITION
\begin{definition}[Power Inverse (PoI)]\index{Power Inverse}
\begin{subequations}
\begin{align}
\left(b^{m}\right)^{\frac{1}{m}} &= b  \label{eq:poi}
\end{align}
\end{subequations}
\end{definition}

%-=-=-= DEFINITION
\begin{definition}[Power Inverse (PoId)]\index{Power Identity}
\begin{subequations}
\begin{align}
1&= b^0  \label{eq:poid1} \\
b^{0}&= 1  \label{eq:poid2}
\end{align}
\end{subequations}
\end{definition}

\begin{notation}[Radical To Power (RTPo)]\index{Powers!Radical to Power}
\begin{align}
\sqrt[m]{b^n} &= b^{\frac{n}{m}}	 \label{eq:rtpo}
\end{align}
\end{notation}

\begin{notation}[Power To Radical (PoTR)]\index{Powers!Power to Radical}
\begin{align}
b^{\frac{n}{m}} &= \sqrt[m]{b^n}	 \label{eq:potr}
\end{align}
\end{notation}

%-=-=-=-=-=-=-=-=-=-=-=-=-=-=-=-=-=-=-=-=-=-=-=-=
%	SUBSECTION: Monomials of Like Terms
%-=-=-=-=-=-=-=-=-=-=-=-=-=-=-=-=-=-=-=-=-=-=-=-=

\subsection{Monomials of Like Terms}\index{Monomials of Like Terms}


%-=-=-= DEFINITION
\begin{definition}[Additive Inverse (AI)]\index{Additive Inverse}
\begin{subequations}
\begin{align}
a + \alert{(-a)} &= 0 \label{eq:ai}
\end{align}
\end{subequations}
\end{definition}

\subsection{Surds}\index{Surds}

%-=-=-= EXAMPLE
\begin{example}[id:20141108-085327] \label{20141108-085327} \index{Example!20141108-085327} \hfill \\

Simplify $2\sqrt{2}-\dfrac{\left(\sqrt{2}\right)^3}{3}-\left(2\left(-\sqrt{2}\right)-\dfrac{\left(-\sqrt{2}\right)^3}{3} \right)$

\soln

\solnsteps
\begin{align*}
2 \sqrt{2}- \dfrac{\alert{1}\left(\alert{1}\sqrt{2}\right)^{3}}{3}-\alert{1} \left(2 \left(-\alert{1}\sqrt{2} \right) - \dfrac{\alert{1} \left(-\alert{1} \sqrt{2}\right)^{3}}{3} \right) && \text{MId} \eqref{eq:mid1} \\
2 \sqrt{2} - \dfrac{1 \left(1\sqrt{2}\right)^{3}}{3}-1 \left(2 \left(\alert{\neg 1} \sqrt{2} \right) - \dfrac{1 \left(\alert{\neg 1} \sqrt{2}\right)^{3}}{3} \right) && \text{ONeg} \eqref{eq:oneg1} \\
2 \sqrt{2} + \dfrac{\neg 1 \left(1\sqrt{2}\right)^{3}}{3}+ \neg 1 \left(2 \left(\alert{\neg 1} \sqrt{2} \right) + \dfrac{\neg 1 \left(\alert{\neg 1} \sqrt{2}\right)^{3}}{3} \right) && \text{DOS} \eqref{eq:dos1} \\
2 \cdot \alert{2^{1/2}} + \dfrac{\neg 1 \left(1 \cdot \alert{2^{1/2}}\right)^{3}}{3} + \neg 1 \left(2 \left(\neg  1 \cdot \alert{2^{1/2}} \right) + \dfrac{\neg 1 \left(\neg  1  \cdot \alert{2^{1/2}}\right)^{3}}{3} \right) && \text{RTPo} \eqref{eq:rtpo} \\
2 \cdot 2^{1/2} + \dfrac{\neg 1 \left(1 \cdot 2^{1/2}\right)^{3}}{3} + \neg 1 \left(\alert{2 \cdot \neg  1 \cdot 2^{1/2}} + \dfrac{\neg 1 \left(\neg 1  \cdot 2^{1/2}\right)^{3}}{3} \right) && \text{JTC} \eqref{eq:jtc} \\
2 \cdot 2^{1/2} + \dfrac{\neg 1 \cdot \alert{1 \cdot 2^{3/2}}}{3} + \neg 1 \left(2 \cdot \neg  1 \cdot 2^{1/2} + \dfrac{\neg 1 \cdot \alert{\neg 1  \cdot 2^{3/2}}}{3} \right) && \text{PoPrPo} \eqref{eq:poprpo1} \\
2 \cdot 2^{1/2} + \dfrac{\neg 1 \cdot 1 \cdot \alert{2^{2/2} \cdot 2^{1/2}}}{3} + \neg 1 \left(2 \cdot \neg  1 \cdot 2^{1/2} + \dfrac{\neg 1 \cdot \neg 1  \cdot \alert{2^{2/2} \cdot 2^{1/2}}}{3} \right) && \text{PrCBPo} \eqref{eq:prcbpo2} \\
2 \cdot 2^{1/2} + \dfrac{\neg 1 \cdot 1 \cdot \alert{2} \cdot 2^{1/2}}{3} + \neg 1 \left(2 \cdot \neg  1 \cdot 2^{1/2} + \dfrac{\neg 1 \cdot \neg 1  \cdot \alert{2} \cdot 2^{1/2}}{3} \right) && \text{MId} \eqref{eq:mid2} \\
2 \cdot \alert{\sqrt{2}} + \dfrac{\neg 1 \cdot 1 \cdot 2 \cdot \alert{\sqrt{2}}}{3} + \neg 1 \left(2 \cdot \neg  1 \cdot \alert{\sqrt{2}} + \dfrac{\neg 1 \cdot \neg 1  \cdot 2 \cdot \alert{\sqrt{2}}}{3} \right) && \text{PoTR} \eqref{eq:potr} \\
2 \cdot \sqrt{2} + \dfrac{\neg 1 \cdot 1 \cdot 2 \cdot \sqrt{2}}{3} + \alert{\neg 1} \cdot 2 \cdot \neg  1 \cdot \sqrt{2} + \dfrac{\alert{\neg 1} \cdot \neg 1 \cdot \neg 1  \cdot 2 \cdot \sqrt{2}}{3}  && \text{DPE} \eqref{eq:dpe1} \\
2 \cdot \sqrt{2} + \dfrac{\neg 2 \cdot \sqrt{2}}{3} + 2 \cdot \sqrt{2} + \dfrac{\neg 2 \cdot \sqrt{2}}{3}  && \text{OOM} \eqref{eq:oom} \\
2\sqrt{2} + \dfrac{\neg 2\sqrt{2}}{3} + 2\sqrt{2} + \dfrac{\neg 2\sqrt{2}}{3} && \text{CTJ} \eqref{eq:ctj} \\
\left(2 + \dfrac{\neg 2}{3} + 2 + \dfrac{\neg 2}{3} \right)  \sqrt{2} && \text{DPF} \eqref{eq:dpf1} \\
\dfrac{8}{3}\sqrt{2} && \text{OOA} \eqref{eq:ooa}
\end{align*}

\qdepend

\qdependlist
example \ref{20141108-083108}-20141108-083108

\end{example}

\end{document}


%\subfile{chapter-myp-univariatepolynomialexpressions}
%\subfile{chapter-myp-multivariatepolynomialexpressions}
%\subfile{chapter-myp-algebraicequations}
%\subfile{chapter-myp-analyticgeometry}
%\part{Base Year}

%\part{DP-SL}
%%!TEX root = /Users/markholson/Dropbox/+Projects/LatexFiles/MathNotebook/20150312-163541-rs2.2N-MarksMathNotebook.tex
%-=-=-=-=-=-=-=-=-=-=-=-=-=-=-=-=-=-=-=-=-=-=-=-=
%
%	CHAPTER 
%
%-=-=-=-=-=-=-=-=-=-=-=-=-=-=-=-=-=-=-=-=-=-=-=-=

\chapterimage{chapter_head_2.pdf} % Chapter heading image

\chapter{Functions}

%-=-=-=-=-=-=-=-=-=-=-=-=-=-=-=-=-=-=-=-=-=-=-=-=
%	SECTION: Evaluating Functions
%-=-=-=-=-=-=-=-=-=-=-=-=-=-=-=-=-=-=-=-=-=-=-=-=

\section{Evaluating Functions}\label{Evaluating Functions}

\begin{example}[id:20141106-083703]\label{20141106-083703}\index{Example!20141106-083703} \hfill \\
Given $f(x)= \dfrac{1}{12}x^3-x^2+4x$, evaluate $f(3)-f(2)$.

\soln

\solnsteps
\begin{align*}
f(3)-f(2) & = \left(\dfrac{\farg{3}^3}{12} - \farg{3}^2 + 4\farg{3} \right) - \left(\dfrac{\farg{2}^3}{12} - \farg{2}^2 + 4\farg{2} \right) \\
&= \left(\dfrac{(3)^3}{12} - (3)^2 + 4(3) \right) - 1 \left(\dfrac{(2)^3}{12} - (2)^2 + 4(2) \right) && \text{MId} \eqref{eq:mid1} \\
&= \left(\dfrac{(3)^3}{12} + \neg  (3)^2 + 4(3)\right) + \neg 1 \left(\dfrac{(2)^3}{12} + \neg (2)^2 + 4(2) \right) && \text{DOS} \eqref{eq:dos1} \\
&= \left(\dfrac{27}{12} + \neg 9 + 4(3) \right)+ \neg 1 \left(\dfrac{8}{12} + \neg 4 + 4(2) \right) && \text{OOE} \eqref{eq:ooe} \\
&= \left(\dfrac{27}{12} + \neg 9 + 12 \right) + \neg 1 \left(\dfrac{8}{12} + \neg 4 + 8 \right) && \text{OOM} \eqref{eq:oom} \\
&= \left(\dfrac{27}{12} + 3 \right) + \neg 1 \left(\dfrac{8}{12} + 4 \right) && \text{OOA} \eqref{eq:ooa} \\
&= \left(\dfrac{27+36}{12} \right) + \neg 1 \left(\dfrac{8+48}{12}\right) && \text{FOOA} \eqref{eq:fooa1} \\
&= \left(\dfrac{63}{12} \right) + \neg 1 \left(\dfrac{56}{12}\right) && \text{OOA} \eqref{eq:ooa} \\
&= \left(\dfrac{63}{12} \right) + \neg \dfrac{56}{12} && \text{OOM} \eqref{eq:oom} \\
&= \dfrac{7}{12} && \text{OOA} \eqref{eq:ooa} 
\end{align*}

\qdepend 

\qdependlist

example \ref{20141105-144223}-20141105-144223


\end{example}

%-=-=-=-=-=-=-=-=-=-=-=-=-=-=-=-=-=-=-=-=-=-=-=-=
%	SECTION: Rational Functions
%-=-=-=-=-=-=-=-=-=-=-=-=-=-=-=-=-=-=-=-=-=-=-=-=

\section{Quadratic Functions}\label{Quadratic Functions}

\subsection{Completing the Square}

\begin{definition}[Completing The Square]
Completing the square is the process used to convert a quadratic polynomial function \(f(x)=ax^2+bx+c\) to the form

\[ f(x)=a \left(x+\dfrac{b}{2a} \right)^2 + \left(c-\dfrac{b^2}{4a} \right) \]

We can simplify this form by defining \(B = \dfrac{b}{2a} \) and \(C=c-\dfrac{b^2}{4a} \), which gives us 

\[ f(x)=a \left(x+B \right) + C \] 

\hfill \cite{mathworld:completethesquare}
\end{definition}

\begin{align*}
f(x) 	&= ax^2+bx+c \\
		&= a \left[x^2+ \dfrac{b}{a}x + \dfrac{c}{a} \right] && \text{DPF} \eqref{eq:dpf2} \\
		&= a \left[x^2 + \dfrac{b}{a}x + \cRed{k} + \cRed{\neg k} + \dfrac{c}{a} \right] && \text{AId} \eqref{eq:aid2} \\
		&= a \left[ \left(x^2 + \dfrac{b}{a}x + \cRed{k} \right)+ \left(\cRed{-k}+ \dfrac{c}{a}  \right) \right] && \text{APA} \eqref{eq:apa2} \\
\end{align*}

\begin{figure}[h!]
\centering
\begin{tikzpicture}[scale=1, auto]

% Place nodes

\node[firstterm](11){$x$}; \node[factoradd,right=of 11](plus1){}; \node[secondterm, right=of plus1](12){$\frac{b}{2a}$};
\node[firstterm, below=of 11](21){$x$}; \node[factoradd,right=of 21](plus2){}; \node[secondterm, right=of plus2](22){$\frac{b}{2a}$};

\node[multiply, below=of 21](31){$x^2$};
\node[multiply, below=of 22](32){$\cRed{k}=\frac{b^2}{4a^2}$};

\node[multiply, right=of 12](13){$\frac{b}{2a}x$};
\node[multiply, right=of 22](23){$\frac{b}{2a}x$};

\node[add, below=of 23](33){$\frac{b}{2}x$};

\path [line](11) edge[bend right=30]node[color=black, midway, left]{$\times$}(21);
\path [line](12) edge[bend left=30]node[color=black,, midway, right]{$\times$}(22);
\path [line](21)--(31);

\path [line](21) edge[bend left=30]node[color=black, pos=0.4, below]{$\times$}(12);
\path [line](11) edge[bend right=30](22);
\path [line](22)--(32);

\path [line](12)--(13);
\path [line](22)--(23);

\path [line](13)--node[color=black, midway, right]{$+$}(23);
\path [line](23)--(33);

\end{tikzpicture}
\caption{Factoring Organizer used to find the value of $\cRed{k}$}
\end{figure}

\begin{align*}
f(x) 	&=  a \left[ \left(x^2 + \dfrac{b}{a}x + \cRed{\dfrac{b^2}{4a^2}} \right)+ \left(\cRed{\dfrac{-b^2}{4a^2}}+ \dfrac{c}{a}  \right) \right]\\
		&= a \left[ \left(x+ \dfrac{b}{2a} \right)^2 + \cRed{\dfrac{-b^2}{4a^2}}+ \dfrac{c}{a} \right]  && \text{DPF} \eqref{eq:dpf2} \\
		&= a \left[  \left(x+ \dfrac{b}{2a} \right)^2 + \dfrac{-b^2}{4a^2}+ \dfrac{4ac}{4a^2} \right] \\ % TODO Common Denominator Step
		&= a \left[  \left(x+ \dfrac{b}{2a} \right)^2 + \dfrac{4ac-b^2}{4a^2} \right]  && \text{OOA} \eqref{eq:ooa} \\
		&= a \left(x+ \dfrac{b}{2a} \right)^2 + \dfrac{4ac-b^2}{4a} && \text{DPE} \eqref{eq:dpe1} \\ 
		&= a \left(x+ \dfrac{b}{2a} \right)^2 + \left(c-\dfrac{b^2}{4a} \right) % TODO Reduce fraction step 
\end{align*}


%-=-=-=-=-=-=-=-=-=-=-=-=-=-=-=-=-=-=-=-=-=-=-=-=
%	SECTION: Trigonometric Functions
%-=-=-=-=-=-=-=-=-=-=-=-=-=-=-=-=-=-=-=-=-=-=-=-=

\section{Trigonometric Functions}\label{Trigonometric Functions}

\subsection{Graphing Trigonometric Function}\label{Graphing Trigonometric Functions}

\newpage
\subsubsection{Graphing the Sine Function}\label{Graphing the Sine Function}

\begin{center}
\begin{tikzpicture}[scale=1, auto]
% Place nodes
\node[task](1){Draw principal axis $y_{sa}=D$ };
\node[decision, right= of 1] (2) {$B<0$};
\node[task, right= of 2] (3) {$(-Bx \pm C)=-1(Bx\mp C)$};
\node[task, below= of 3] (4) {$A \sin [-1(Bx\mp C)]= -A \sin(Bx \mp C)$};
\node[task, left= of 4] (5) {$P= \dfrac{360 \degree}{ \lvert B \rvert}$};
\node[task, left= of 5] (6) {$PS = -\dfrac{C}{B}$};

\node[abscissa, below= of 5, node distance=2.0cm] (7) {$x_0= PS$};
\node[abscissa, below= of 7] (8) {$x_1 = PS + \dfrac{P}{4}$};
\node[abscissa, below= of 8] (9) {$x_2=PS + \dfrac{P}{2}$};
\node[abscissa, below= of 9] (10) {$x_3= PS + \dfrac{3P}{4}$};
\node[abscissa, below= of 10] (11) {$x_4= PS + P$};

\node[task, below= of 11] (12) {$ a=\lvert A \rvert$};
\node[decision, below= of 12] (13) {$A>0$};

\node[ordinate, left= of 11] (14) {$y_4=D$};
\node[ordinate, above= of 14] (15) {$y_3=D-a$};
\node[ordinate, above= of 15] (16) {$y_2=D$};
\node[ordinate, above= of 16] (17) {$y_1=D+a$};
\node[ordinate, above= of 17] (18) {$y_0=D$};

\node[ordinate, right= of 11] (19) {$y_4=D$};
\node[ordinate, above= of 19] (20) {$y_3=D+a$};
\node[ordinate, above= of 20] (21) {$y_2=D$};
\node[ordinate, above= of 21] (22) {$y_1=D-a$};
\node[ordinate, above= of 22] (23) {$y_0=D$};

\node[bluedotted, fit =(7)(8)(9)(10)(11)](boxabscissa) {};
\node [above of =7,  node distance=1.2cm] {\textcolor{sthlmBlue}{\textbf{Abscissae}}};

\node[reddotted, fit =(14)(15)(16)(17)(18)](boxordinate1) {};
\node [above of =23,  node distance=1.2cm] {\textcolor{sthlmRed}{\textbf{Ordinates}}};

\node[reddotted, fit =(19)(20)(21)(22)(23)](boxordinate2) {};
\node [above of =18,  node distance=1.2cm] {\textcolor{sthlmRed}{\textbf{Ordinates}}};

\node[pointcircle, left=of 14](25) {$Q_4$};
\node[pointcircle, left=of 15](26) {$Q_3$};
\node[pointcircle, left=of 16](27) {$Q_2$};
\node[pointcircle, left=of 17](28) {$Q_1$};
\node[pointcircle, left=of 18](29) {$Q_0$};

\node[pointcircle, right=of 19](30) {$Q_4$};
\node[pointcircle, right=of 20](31) {$Q_3$};
\node[pointcircle, right=of 21](32) {$Q_2$};
\node[pointcircle, right=of 22](33) {$Q_1$};
\node[pointcircle, right=of 23](34) {$Q_0$};
% Draw Edges
\path [line](1)--(2);
\path [line](2)-- node[color=black]{Yes}(3);
\path [line](3)--(4);
\path [line](4)--(5);
\path [line](2)--node[color=black]{No}(5);
\path [line](5)--(6);
\path [line](6) edge   [bend left=10] (7);
\path [line](7)--(8);
\path [line](8)--(9);
\path [line](9)--(10);
\path [line](10)--(11);
\path [line](12)--(13);
\path [line](11)--(12);
\path [line](13) edge [bend left=30] node[color=black, left]{Yes}(14);
\path [line](13) edge [bend left=-30] node[color=black, right]{No}(19);
\path [line](14)--(15);
\path [line](15)--(16);
\path [line](16)--(17);
\path [line](17)--(18);
\path [line](19)--(20);
\path [line](20)--(21);
\path [line](21)--(22);
\path [line](22)--(23);
\end{tikzpicture}
\end{center}

%-=-=-= EXAMPLE
\begin{example}[id:20150327-115953] \label{20150327-115953}\index{Example!20150327-115953} \hfill \\
Sketch the function $g(x)=\sin(x)+2$

\soln

\begin{sagesilent}
a20150327115953=1
b20150327115953=1
c20150327115953=0
d20150327115953=2

fbase20150327115953=sin(x*pi/180)

f20150327115953=a20150327115953*sin(b20150327115953*(x*pi/180)+c20150327115953*(pi/180))+d20150327115953
ha20150327115953=d20150327115953

phaseshift20150327115953=-c20150327115953/b20150327115953
amplitude20150327115953=abs(a20150327115953)
period20150327115953=360/abs(b20150327115953)
fmaxg20150327115953=d20150327115953+abs(a20150327115953)+0.5
fming20150327115953=d20150327115953-abs(a20150327115953)-0.5
fmax20150327115953=d20150327115953+abs(a20150327115953)
fmin20150327115953=d20150327115953-abs(a20150327115953)

x020150327115953=phaseshift20150327115953
x120150327115953=phaseshift20150327115953+(period20150327115953/4)
x220150327115953=x120150327115953+(period20150327115953/4)
x320150327115953=x220150327115953+(period20150327115953/4)
x420150327115953=x320150327115953+(period20150327115953/4)

x0plot20150327115953=line([(x020150327115953,fmax20150327115953),(x020150327115953,fmin20150327115953)], linestyle="--",zorder=3, rgbcolor=(0.364705882352941, 0.137254901960784, 0.490196078431373), thickness=2)
x1plot20150327115953=line([(x120150327115953,fmax20150327115953),(x120150327115953,fmin20150327115953)], linestyle="--",zorder=3, rgbcolor=(0.364705882352941, 0.137254901960784, 0.490196078431373), thickness=2)
x2plot20150327115953=line([(x220150327115953,fmax20150327115953),(x220150327115953,fmin20150327115953)], linestyle="--",zorder=3, rgbcolor=(0.364705882352941, 0.137254901960784, 0.490196078431373), thickness=2)
x3plot20150327115953=line([(x320150327115953,fmax20150327115953),(x320150327115953,fmin20150327115953)], linestyle="--",zorder=3, rgbcolor=(0.364705882352941, 0.137254901960784, 0.490196078431373), thickness=2)
x4plot20150327115953=line([(x420150327115953,fmax20150327115953),(x420150327115953,fmin20150327115953)], linestyle="--",zorder=3, rgbcolor=(0.364705882352941, 0.137254901960784, 0.490196078431373), thickness=2)

point020150327115953=point((x020150327115953,f20150327115953(x=x020150327115953)), rgbcolor=(0.364705882352941, 0.137254901960784, 0.490196078431373), size=50)
point120150327115953=point((x120150327115953,f20150327115953(x=x120150327115953)), rgbcolor=(0.364705882352941, 0.137254901960784, 0.490196078431373), size=50)
point220150327115953=point((x220150327115953,f20150327115953(x=x220150327115953)), rgbcolor=(0.364705882352941, 0.137254901960784, 0.490196078431373), size=50)
point320150327115953=point((x320150327115953,f20150327115953(x=x320150327115953)), rgbcolor=(0.364705882352941, 0.137254901960784, 0.490196078431373), size=50)
point420150327115953=point((x420150327115953,f20150327115953(x=x420150327115953)), rgbcolor=(0.364705882352941, 0.137254901960784, 0.490196078431373), size=50)

if a20150327115953>0:
    if b20150327115953>0:
        maxmin20150327115953="not reflected"
    else:
        maxmin20150327115953="reflected"
else:
    if b20150327115953>0:
        maxmin20150327115953="reflected"
    else:
        maxmin20150327115953="not reflected"
    
#Define the plots for the trigonometric and sinusoidal axis functions
plotf120150327115953=plot(f20150327115953,(x,x020150327115953,x420150327115953),rgbcolor=(0, 0.431372549019608, 0.749019607843137),thickness=3)
plotf220150327115953=plot(f20150327115953,(x,-180,365),rgbcolor=(0, 0.431372549019608, 0.749019607843137),thickness=3)
plotha20150327115953=plot(ha20150327115953,(x,-180,365),rgbcolor=(0.768627450980392, 0, 0.392156862745098),thickness=3)
plotfbase20150327115953=plot(fbase20150327115953,(x,-180,365), ymin=-1,rgbcolor=(0.2, 0.2, 0.2),thickness=2)

#Define the plots to include in the final plot
plotx20150327115953=plotha20150327115953+x0plot20150327115953+x1plot20150327115953+x2plot20150327115953+x3plot20150327115953+x4plot20150327115953
plotspts20150327115953=plotha20150327115953+x0plot20150327115953+x1plot20150327115953+x2plot20150327115953+x3plot20150327115953+x4plot20150327115953+point020150327115953+point120150327115953+point220150327115953+point320150327115953+point420150327115953
plotfint20150327115953=plotf120150327115953+plotha20150327115953+x0plot20150327115953+x1plot20150327115953+x2plot20150327115953+x3plot20150327115953+x4plot20150327115953+point020150327115953+point120150327115953+point220150327115953+point320150327115953+point420150327115953
plotfcompare20150327115953=plotf220150327115953+plotha20150327115953+plotfbase20150327115953

\end{sagesilent}

\solnsteps

\begin{align*}
	f(x)&=A \sin(Bx+C)+D\\
	g(x)&= \sin(x)+ 2\\
	g(x)&= 1 \sin(1x)+2 && \text{MId} \eqref{eq:mid1} \\
	g(x)&= 1 \sin(1x+0)+2 && \text{AId} \eqref{eq:aid1} 
\end{align*}

Thus, A=1, B=1, C=0 and D=2\\

\textbf{Step 1:}  Find and draw the sinusoidal-axis. \\

\begin{align*}
	y_{sa} &= D \\
	y_{sa} &= 2
\end{align*}


$\sageplot{plotha20150327115953,ymin=fming20150327115953,ymax=fmaxg20150327115953,axes_labels=['$x$','$y$'],ticks=[[-180,-150,-120,-90,-60,-30,0,30, 60,90,120,150,180,210,240,270,300,330,360],1],figsize=[5.5,2.5],tick_formatter="latex",fontsize=7, gridlines="True", frame="True",transparent="True"}$ 

\textbf{Step 2:}  Check if the value of $B<0$ \\

Since $B>0$, there is nothing to do here. \\

\textbf{Step 3:}  Calculate the period $P$

\begin{align*}
	P &= \dfrac{360 \degree}{\abs{B}} \\
	P &= \dfrac{360 \degree}{\abs{\farg{1}}} && \text{SPE} \eqref{eq:spe} \\
	P & = 360 \degree
\end{align*}

\textbf{Step 4:}  Calculate the phase shift $PS$


\begin{align*}
Bx_{ps}+C &=0 \\
1x_{px}+0 &=0 && \text{SPE} \eqref{eq:spe} \\
1x_{px} &= 0 && \text{AId} \eqref{eq:aid2} \\
x_{ps} &=0 && \text{MId} \eqref{eq:mid2} 
\end{align*}

\textbf{Step 5:}  State the abscissae of the some points of interest.

\begin{align*}
x_0 &= PS \\
x_0 &= 0 \degree \\ 
& \\
x_1 &= PS + \dfrac{P}{4} \\
x_1 &= 0 \degree + \dfrac{\farg{360 \degree}}{4} && \text{SPE} \eqref{eq:spe} \\
x_1 &= 0 \degree + 90 \degree \\ %TODO Reduce Fraction
x_1 &= 90 \degree \\ 
& \\
x_2 &= PS + \dfrac{P}{2} \\
x_2 &= 0 \degree + \dfrac{\farg{360 \degree}}{2} && \text{SPE} \eqref{eq:spe} \\
x_2 &= 0 \degree + 180 \degree \\ %TODO Reduce Fraction
x_2 &= 180 \degree \\ 
& \\
x_3 &= PS + \dfrac{3P}{4} \\
x_3 &= 0 \degree + \dfrac{3 \farg{360 \degree}}{4} && \text{SPE} \eqref{eq:spe} \\
x_3 &= 0 \degree + 270 \degree \\ %TODO Reduce Fraction
x_3 &= 270 \degree \\ 
& \\
x_4 &= PS + P \\
x_4 &= 0 \degree + \farg{360 \degree} && \text{SPE} \eqref{eq:spe} \\
x_4 &= 0 \degree + 360 \degree \\ 
x_4 &= 360 \degree
\end{align*}

\textbf{Step 6:}  Sketch the vertical equations of the lines to indicate the $x$ values of the points that we will plot.


$\sageplot{plotx20150327115953,ymin=fming20150327115953,ymax=fmaxg20150327115953,axes_labels=['$x$','$y$'],ticks=[[-180,-150,-120,-90,-60,-30,0,30, 60,90,120,150,180,210,240,270,300,330,360],1],figsize=[5.5,2.5],tick_formatter="latex",fontsize=7, gridlines="True", frame="True",transparent="True"}$ 


\textbf{Step 7:}  Find the amplitude of the function.

\begin{align*}
a & = \abs{A} \\
a & = \abs{\farg{1}} \\
a &= 1 
\end{align*}

\textbf{Step 8:}  Determine if the functions $y$ values are reflected about the $x$-axis.\\

Since the value of $A=1$ therefore $A>0$ and consequently there will be no reflection of the $y$-values about the $x$-axis. \\

\textbf{Step 9:}  Find the corresponding ordinates to the abscissae fond in step 5.

\begin{align*}
y_4 & = D \\
y_4 &= \farg{2} \\
y_4 &= 2 \\
& \\
y_3 & = D-a \\
y_3 &= \farg{2}-\farg{1} \\
y_3 &= 1 \\
& \\
y_2 & = D \\
y_2 &= \farg{2} \\
y_2 &= 2 \\
& \\
y_1 & = D+a \\
y_1 &= \farg{2}+\farg{1} \\
y_1 &= 3 \\
& \\
y_0 & = D \\
y_0 &= \farg{2} \\
y_0 &= 2
\end{align*}

\textbf{Step 10:}  Plot the points of interest.\\

$\sageplot{plotspts20150327115953,ymin=fming20150327115953,ymax=fmaxg20150327115953,axes_labels=['$x$','$y$'],ticks=[[-180,-150,-120,-90,-60,-30,0,30, 60,90,120,150,180,210,240,270,300,330,360],1],figsize=[5.5,2.5],tick_formatter="latex",fontsize=7, gridlines="True", frame="True",transparent="True"}$ 


\textbf{Step 10:}  Sketch the function.\\

$\sageplot{plotfint20150327115953,ymin=fming20150327115953,ymax=fmaxg20150327115953,axes_labels=['$x$','$y$'],ticks=[[-180,-150,-120,-90,-60,-30,0,30, 60,90,120,150,180,210,240,270,300,330,360],1],figsize=[5.5,2.5],tick_formatter="latex",fontsize=7, gridlines="True", frame="True",transparent="True"}$ 

$\sageplot{plotfcompare20150327115953,axes_labels=['$x$','$y$'],ticks=[[-180,-150,-120,-90,-60,-30,0,30, 60,90,120,150,180,210,240,270,300,330,360],1],figsize=[5.5,2.5],tick_formatter="latex",fontsize=7, gridlines="True", frame="True",transparent="True"}$ 

\end{example}



\newpage
\subsubsection{Graphing the Cosine Function}\label{Graphing the Cosine Function}

\begin{center}
\begin{tikzpicture}[scale=1, auto]
% Place nodes
\node[task](1){Draw principal axis $y_{sa}=D$ };
\node[decision, right= of 1] (2) {$B<0$};
\node[task, right= of 2] (3) {$(-Bx\pm C)=-1(Bx \mp C)$};
\node[task, below= of 3] (4) {$A \cos [-1(Bx \mp C)]=A \cos(Bx \mp C)$};
\node[task, left= of 4] (5) {$P= \dfrac{360 \degree}{ \lvert B \rvert}$};
\node[task, left= of 5] (6) {$PS = -\dfrac{C}{B}$};

\node[abscissa, below= of 5, node distance=2.0cm] (7) {$x_0= PS$};
\node[abscissa, below= of 7] (8) {$x_1 = PS + \dfrac{P}{4}$};
\node[abscissa, below= of 8] (9) {$x_2=PS + \dfrac{P}{2}$};
\node[abscissa, below= of 9] (10) {$x_3= PS + \dfrac{3P}{4}$};
\node[abscissa, below= of 10] (11) {$x_4= PS + P$};

\node[task, below= of 11] (12) {$ a=\lvert A \rvert$};
\node[decision, below= of 12] (13) {$A>0$};

\node[ordinate, left= of 11] (14) {$y_4=D+a$};
\node[ordinate, above= of 14] (15) {$y_3=D$};
\node[ordinate, above= of 15] (16) {$y_2=D-a$};
\node[ordinate, above= of 16] (17) {$y_1=D$};
\node[ordinate, above= of 17] (18) {$y_0=D+a$};

\node[ordinate, right= of 11] (19) {$y_4=D-a$};
\node[ordinate, above= of 19] (20) {$y_3=D$};
\node[ordinate, above= of 20] (21) {$y_2=D+a$};
\node[ordinate, above= of 21] (22) {$y_1=D$};
\node[ordinate, above= of 22] (23) {$y_0=D-a$};

\node[bluedotted, fit =(7)(8)(9)(10)(11)](boxabscissa) {};
\node [above of =7,  node distance=1.2cm] {\textcolor{sthlmBlue}{\textbf{Abscissae}}};

\node[reddotted, fit =(14)(15)(16)(17)(18)](boxordinate1) {};
\node [above of =23,  node distance=1.2cm] {\textcolor{sthlmRed}{\textbf{Ordinates}}};

\node[reddotted, fit =(19)(20)(21)(22)(23)](boxordinate2) {};
\node [above of =18,  node distance=1.2cm] {\textcolor{sthlmRed}{\textbf{Ordinates}}};
% Draw Edges
\path [line](1)--(2);
\path [line](2)-- node[color=black]{Yes}(3);
\path [line](3)--(4);
\path [line](4)--(5);
\path [line](2)--node[color=black]{No}(5);
\path [line](5)--(6);
\path [line](6) edge   [bend left=10] (7);
\path [line](7)--(8);
\path [line](8)--(9);
\path [line](9)--(10);
\path [line](10)--(11);
\path [line](12)--(13);
\path [line](11)--(12);
\path [line](13) edge [bend left=30] node[color=black, left]{Yes}(14);
\path [line](13) edge [bend left=-30] node[color=black, right]{No}(19);
\path [line](14)--(15);
\path [line](15)--(16);
\path [line](16)--(17);
\path [line](17)--(18);
\path [line](19)--(20);
\path [line](20)--(21);
\path [line](21)--(22);
\path [line](22)--(23);

\node[pointcircle, left=of 14](25) {$Q_4$};
\node[pointcircle, left=of 15](26) {$Q_3$};
\node[pointcircle, left=of 16](27) {$Q_2$};
\node[pointcircle, left=of 17](28) {$Q_1$};
\node[pointcircle, left=of 18](29) {$Q_0$};

\node[pointcircle, right=of 19](30) {$Q_4$};
\node[pointcircle, right=of 20](31) {$Q_3$};
\node[pointcircle, right=of 21](32) {$Q_2$};
\node[pointcircle, right=of 22](33) {$Q_1$};
\node[pointcircle, right=of 23](34) {$Q_0$};
\end{tikzpicture}
\end{center}


%-=-=-=-=-=-=-=-=-=-=-=-=-=-=-=-=-=-=-=-=-=-=-=-=
%	SECTION: Rational Functions
%-=-=-=-=-=-=-=-=-=-=-=-=-=-=-=-=-=-=-=-=-=-=-=-=

\section{Rational Functions}\label{Rational Functions}

%-=-=-= DEFINITION
\begin{definition}[Asympote]\index{Asympote}

An asymptote is a line or curve that apporoaches a given curve arbitrarily close.\cite{mathworld:asymptote} \\

A vertical asymptote is a vertical line $x_{va}=c$, that is used to visualize the values of $x$ for which the function is not defined. 

\end{definition}

\begin{figure}[ht]
\begin{center}
	\begin{tikzpicture}
	\begin{axis}[
            domain=-6:2,
            ymax=4,
            ymin=-4,
            samples=100,
            axis lines =middle, xlabel=$x$, ylabel=$y$,
            every axis y label/.style={at=(current axis.above origin),anchor=south},
            every axis x label/.style={at=(current axis.right of origin),anchor=west},
			restrict y to domain=-20:20
          ]
          \addplot [dashed, stockholmPink, smooth] plot coordinates {(-2,4) (-2,-4)}; %% {.451};

          \addplot [very thick, stockholmBlue, smooth] {1/(x+2)};

          \node at (axis cs:-5.7,3.8) [anchor=west] {\color{stockholmPink}Vertical Asymptote};  

        \end{axis}
\end{tikzpicture}

\end{center}
\caption{Vertical Asymptote}
\label{figure:rectangularhyperbola}
\end{figure}

%-=-=-= EXAMPLE
\begin{example}[id:20141111-190212] \label{20141111-190212}\index{Example!20141111-190212} \hfill \\

Find the vertical asymptote of the function $R(x)=\dfrac{7}{x+8}$

\soln

\solnsteps

We are interested in the values of $x$ for which the denominator of $R(x)$ has a value of zero.

\begin{align*}
x_{va}+8 &= 0 \\
x_{va} &=-8  &&\text{solving} 
\end{align*}

\begin{center}
	\begin{tikzpicture}
	\begin{axis}[
            domain=-20:10,
            ymax=10,
            ymin=-10,
            samples=100,
            axis lines =middle, xlabel=$x$, ylabel=$y$,
            every axis y label/.style={at=(current axis.above origin),anchor=south},
            every axis x label/.style={at=(current axis.right of origin),anchor=west},
			restrict y to domain=-20:20
          ]
          \addplot [dashed, stockholmPink, smooth] plot coordinates {(-8,10) (-8,-10)}; %% {.451};

          \addplot [very thick, stockholmBlue, smooth] {7/(x+8)};

          \node at (axis cs:-17.7,5.8) [anchor=west] {\color{stockholmPink}$x_{va}=-8$};  

        \end{axis}
\end{tikzpicture}

\end{center}
\end{example}

%-=-=-= EXAMPLE
\begin{example}[id:20141111-192213] \label{20141111-192213}\index{Example!20141111-192213} \hfill \\

Find the vertical asymptote of the function $R(x)=\dfrac{x+4}{2x+5}$

\soln

\solnsteps

We are interested in the values of $x$ for which the denominator of $R(x)$ has a value of zero.

\begin{align*}
2x_{va}+5 &= 0 \\
x_{va} &= -\dfrac{5}{2}  &&\text{solving \ref{20141111-215726}} \\
\end{align*}

\begin{center}
	\begin{tikzpicture}
	\begin{axis}[
            domain=-6:2,
            ymax=5,
            ymin=-5,
            samples=100,
            axis lines =middle, xlabel=$x$, ylabel=$y$,
            every axis y label/.style={at=(current axis.above origin),anchor=south},
            every axis x label/.style={at=(current axis.right of origin),anchor=west},
			restrict y to domain=-5:5
          ]
          \addplot [dashed, stockholmPink, smooth] plot coordinates {(-5/2,4) (-5/2,-4)}; %% {.451};

          \addplot [very thick, stockholmBlue, smooth] {(x+4)/(2*x+5)};

          \node at (axis cs:-4.7,2.8) [anchor=west] {\color{stockholmPink}$x_{va}=-\frac{5}{2}$};  

        \end{axis}
\end{tikzpicture}

\end{center}
\end{example}








%%!TEX root = /Users/markholson/Dropbox/+Projects/LatexFiles/MathNotebook/20150312-163541-rs2.2N-MarksMathNotebook
%-=-=-=-=-=-=-=-=-=-=-=-=-=-=-=-=-=-=-=-=-=-=-=-=
%
%	CHAPTER 
%
%-=-=-=-=-=-=-=-=-=-=-=-=-=-=-=-=-=-=-=-=-=-=-=-=

\chapterimage{chapter_head_2.pdf} % Chapter heading image

\chapter{Vectors}

%-=-=-=-=-=-=-=-=-=-=-=-=-=-=-=-=-=-=-=-=-=-=-=-=
%	SECTION: 
%-=-=-=-=-=-=-=-=-=-=-=-=-=-=-=-=-=-=-=-=-=-=-=-=
\section{Component \& Geometric Vectors}


\section{Dot Product}

%-=-=-= DEFINITION
\begin{definition}[Dot Product in $\mathbb{R}^2$ (VDPr)]\index{Dot Product $\mathbb{R}^2$}
The dot product of vectors $\vec{u}=(u_x,u_y)$ and $\vec{v}=(v_x, v_y)$ is 

\begin{align}
\vec{u} \cdot \vec{v} 	&= u_x v_x + u_y v_y  \label{eq:vdpr1}
\end{align}

The dot product of vectors $\vec{u}=u_x \ihat + u_y \jhat$ and $\vec{v}=v_x \ihat + v_y \jhat$ is 

\begin{align}
\vec{u} \cdot \vec{v} 	&= u_xv_x \ihat + u_yv_y \jhat \label{eq:vdpr2}
\end{align}

\end{definition}

%-=-=-= EXAMPLE
\begin{example}[id:20141214-195907] \label{20141214-195907}\index{Example!20141214-195907} \hfill \\

Find the dot product of the vectors $\vec{u}=(-5,12)$ and $\vec{v}=(-7,- 2)$

\soln

\solnsteps
\begin{align*}
\vec{u} \cdot \vec{v} &= u_x v_x + u_y v_y && \text{VDPr} \eqref{eq:vdpr1} \\
&=\farg{-5}\farg{-7} + \farg{12}\farg{-2} && \text{SPE} \eqref{eq:spe} \\
&= 35 + -24  && \text{OOM} \eqref{eq:oom} \\
&= 11
\end{align*}
\end{example}

%-=-=-= EXAMPLE
\begin{example}[id:20141214-194701] \label{20141214-194701}\index{Example!20141214-194701} \hfill \\

Find the dot product of the vectors $\vec{u}=-10\ihat - 10 \jhat$ and $\vec{v}=-9\ihat + 7 \jhat$

\soln
\solnsteps
\begin{align*}
\vec{u} \cdot \vec{v} & = \cBlue{u_x}\cRed{v_x} \ihat + \cBlue{u_y}\cRed{v_y} \jhat && \text{VDPr} \eqref{eq:vdpr2} \\
&= \farg{-10}\farg{-9}+ \farg{-10}\farg{7} && \text{SPE} \eqref{eq:spe} \\
&= 90+\neg 70 && \text{OOM} \eqref{eq:oom} \\
&= 20 && \text{OOA} \eqref{eq:ooa} 
\end{align*}


\end{example}
%\subfile{chapter-dpsl-differentiation}
%\documentclass[20150903-160354-rs2.2-MarksMathNotebook.tex]{subfiles}

\begin{document}
%-=-=-=-=-=-=-=-=-=-=-=-=-=-=-=-=-=-=-=-=-=-=-=-=
%
%	CHAPTER
%
%-=-=-=-=-=-=-=-=-=-=-=-=-=-=-=-=-=-=-=-=-=-=-=-=

\chapterimage{Pictures/chapter_head_2.pdf} % Chapter heading image

\chapter{Differentiation}

%-=-=-=-=-=-=-=-=-=-=-=-=-=-=-=-=-=-=-=-=-=-=-=-=
%	SECTION:
%-=-=-=-=-=-=-=-=-=-=-=-=-=-=-=-=-=-=-=-=-=-=-=-=
\section{Limit of the Difference Quotient}

%-=-=-= DEFINITION
\begin{definition}[Derivative]\index{Derivative}

The derivative of a function $f(x)$ with respect to the variable $x$ is defined as
\begin{align}
f'(\cBlue{x}) & \equiv \displaystyle \lim_{\cRed{\Delta x} \to 0} \underbrace{\dfrac{f(\cBlue{x}+\cRed{\Delta x})-f(\cBlue{x})}{\cRed{\Delta x}}}_{\text{Difference Quotient}} \label{eq:dbfp1}
\end{align}

\end{definition}
\begin{figure}[h!]
\centering
\begin{tikzpicture}
	\begin{axis}[
            domain=0:2, range=0:6,ymax=6,ymin=0,
            axis lines =left, xlabel=$x$, ylabel=$y$,
            every axis y label/.style={at=(current axis.above origin),anchor=south},
            every axis x label/.style={at=(current axis.right of origin),anchor=west},
            xtick={1,1.666}, ytick={1,3},
            xticklabels={$x$,$x+\Delta x$}, yticklabels={$f(x)$,$f(x+\Delta x)$},
            axis on top,
          ]
          \addplot [sthlmRed!15, domain=(0:2)] {-3.348+4.348*x};
          \addplot [sthlmRed!32, domain=(0:2)] {-2.704+3.704*x};
          \addplot [sthlmRed!49, domain=(0:2)] {-1.994+2.994*x};
          \addplot [sthlmRed!66, domain=(0:2)] {-1.326+2.326*x};
          \addplot [sthlmRed!83, domain=(0:2)] {-0.666+1.666*x};
	  \addplot [sthlmDarkGrey,dashed] plot coordinates {(1,0) (1,1)};
          \addplot [sthlmDarkGrey,dashed] plot coordinates {(0,1) (1,1)};
          \addplot [sthlmDarkGrey,dashed] plot coordinates {(0,3) (1.666,3)};
          \addplot [sthlmDarkGrey,dashed] plot coordinates {(1.666,0) (1.666,3)};
          \addplot [very thick,sthlmBlue, smooth,domain=(0:1.833)] {-1/(x-2)};
          \addplot[color=sthlmDarkGrey,fill=sthlmDarkGrey,only marks,mark=*] coordinates{(1.666,3)};  %% closed hole
          \addplot[color=sthlmDarkGrey,fill=sthlmDarkGrey,only marks,mark=*] coordinates{(1,1)};  %% closed hole
          \addplot [very thick,sthlmRed, smooth,domain=(0:2)] {x};
        \end{axis}
\end{tikzpicture}
\caption{\cite{mooculus:textbook}}
\end{figure}
%-=-=-= EXAMPLE
\begin{example}[id:20141219-212546] \label{20141219-212546}\index{Example!20141219-212546} \hfill \\

Differentiate the function $f(x)=5$

\soln

\solnsteps
\begin{align*}
f(x) & = 5x^0 && \text{PoID} \eqref{eq:poid1} \\
f'(x) &= \displaystyle \lim_{\Delta x \to 0} \dfrac{5\farg{x+\Delta x}^0-5\farg{x}^0}{\Delta x} && \text{SPE} \eqref{eq:spe} \& \text{DBFP} \eqref{eq:dbfp1}\\
f'(x) &= \displaystyle \lim_{\Delta x \to 0} \dfrac {5(1)-5(1)}{\Delta x} && \text{PoID} \eqref{eq:poid2} \\
f'(x) &= \displaystyle \lim_{\Delta x \to 0} 0 && \text{OOM} \eqref{eq:oom} \\
f'(x) &= 0 % TODO Evaluate the limit
\end{align*}

\end{example}

%-=-=-=-=-=-=-=-=-=-=-=-=-=-=-=-=-=-=-=-=-=-=-=-=
%	SECTION:
%-=-=-=-=-=-=-=-=-=-=-=-=-=-=-=-=-=-=-=-=-=-=-=-=
\section{Derivative of a Monomial Functions}

\begin{arule}[Derivative of a Constant (DC)]\index{Derivatives! Derivative of a Constant}
\begin{align}
	\left[c \right]' &= 0 \label{eq:dc1} \\
	\dfrac{\dd}{\dx} \left[c \right]&=0  \label{eq:dc2}
\end{align}
\end{arule}

\begin{arule}[Derivative of a Constant Multiple (DCM)]\index{Derivatives! Derivative of a Constant Multiple}
\begin{align}
	\left[c f(x) \right]' &= c \left[f(x)\right]' \label{eq:dcm1} \\
	\dfrac{\dd}{\dx}\left[c f(x) \right] &= c \dfrac{\dd}{\dx}\left[f(x)\right]  \label{eq:dcm2}
\end{align}
\end{arule}

\begin{arule}[Derivative of a Power (DPo)]\index{Derivatives! Derivative of a Power}
\begin{align}
	\left[x^n \right]'&= n x^{n-1} \label{eq:dpo1} \\
	\dfrac{\dd}{\dx} \left[x^n \right] &= n x^{n-1}  \label{eq:dpo2}
\end{align}
\end{arule}

%-=-=-= EXAMPLE
\begin{example}[id:20141124-153017] \label{20141124-153017}\index{Example!20141124-153017} \hfill \\

Differentiate $f(x)=-3$

\soln

\solnsteps
\begin{align*}
f'(x) &= [-3]' && \text{SPE} \eqref{eq:spe} \\
f'(x) &= 0 && \text{DC} \eqref{eq:dc1}
\end{align*}

\qdepend

\qdependlist

example \ref{20141124-152503}-20141124-152503


\end{example}


%-=-=-= EXAMPLE
\begin{example}[id:20141124-141850] \label{20141124-141850}\index{Example!20141124-141850} \hfill \\

Differentiate $f(x)=x^2$

\soln

\solnsteps
\begin{align*}
f'(x) &= \left[x^2\right]' && \text{SPE} \eqref{eq:spe} \\
f'(x) &= \cRed{2}x^{\cRed{2-1}} && \text{DPo} \eqref{eq:dpo1} \\
f'(x) &= 2x^1 && \text{OOA} \eqref{eq:ooa} \\
f'(x) &= 2x && \text{MId} \eqref{eq:mid2}
\end{align*}

\soln

\lesssteps
\begin{align*}
f'(x) & = 2x && \text{DPo} \eqref{eq:dpo1}
\end{align*}

\qdepend

\qdependlist

example \ref{20141124-152503}-20141124-152503



\end{example}

%-=-=-=-=-=-=-=-=-=-=-=-=-=-=-=-=-=-=-=-=-=-=-=-=
%	SECTION:
%-=-=-=-=-=-=-=-=-=-=-=-=-=-=-=-=-=-=-=-=-=-=-=-=
\section{Derivative of Polynomial Functions}

\begin{arule}[Derivative of a Sum (DS)]\index{Derivatives! Derivative of a Sum}
\begin{align}
	[f(x) + g(x)]' &= f'(x) + g'(x) \label{eq:ds1} \\
	\dfrac{\dd}{\dx} \left[f(x) + g(x) \right] &= \dfrac{\dd}{\dx} \left[f(x) \right] + \dfrac{\dd}{\dx} \left[g(x) \right]  \label{eq:ds2}
\end{align}
\end{arule}

%-=-=-= EXAMPLE
\begin{example}[id:20141124-152503] \label{20141124-152503}\index{Example!20141124-152503} \hfill \\

Differentiate $f(x)=x^2-3$

\soln

\solnsteps
\begin{align*}
f(x) &= x^2 + \neg 3 && \text{DOS} \eqref{eq:dos1} \\
f'(x) &= \farg{x^2 + \neg 3}' && \text{SPE} \eqref{eq:spe} \\
f'(x) &= \farg{x^2}' + \farg{\neg 3}' && \text{DS} \eqref{eq:ds1} \\
f'(x) &= \farg{x^2}' + 0 && \text{DC} \eqref{eq:dc1} \\
f'(x) &= \farg{x^2}' && \text{AId} \eqref{eq:aid1} \\
f'(x) &= 2x && \text{DPo} \eqref{eq:dpo1}  \text{\, goto \,}  \ref{20141124-141850}
\end{align*}

\soln

\lesssteps

\begin{align*}
f(x) & = 2x^2 && \text{DPo} \eqref{eq:dpo1} \& \text{DC} \eqref{eq:dpo1}
\end{align*}

\qdepend

\qdependlist

example \ref{20141124-205219}-20141124-205219

\end{example}

%-=-=-= EXAMPLE
\begin{example}[id:20141128-151834] \label{20141128-151834}\index{Example!20141128-151834} \hfill \\

Differentiate $f(x)=3x^2-6x+4$

\soln

\solnsteps
\begin{align*}
f(x) &= 3x^2 + \neg 6x + 4 && \text{DOS} \eqref{eq:dos1} \\
f'(x) &= \farg{3x^2 + \neg 6x + 4}' && \text{SPE} \eqref{eq:spe} \\
f'(x) &= \farg{3x^2}' + \farg{6x}' + \farg{4}' && \text{DS} \eqref{eq:ds1} \\
f'(x) &= \farg{3x^2}' + \farg{6x}' + 0 && \text{DC} \eqref{eq:dc1} \\
f'(x) &= \farg{3x^2}' + \farg{6x}' && \text{AId} \eqref{eq:aid1} \\
f'(x) &= 3\farg{x^2}' + 6\farg{x}' && \text{DCM} \eqref{eq:dcm1} \\
f'(x) &= 3 \left(2x \right) + 6 \left(1\right) && \text{DPo} \eqref{eq:dpo1} \\
f'(x) &= 6x+6 && \text{OOM} \eqref{eq:oom}
\end{align*}

\soln
\lesssteps
\begin{align*}
f'(x) & = 6x+6 && \text{DS} \eqref{eq:ds1}
\end{align*}
\end{example}

\subsection{Finding the vertex of a quadratic function using differentiation.}

We can find the vertex of a quadratic function, $f(x)$ using differentiation by:

\begin{enumerate}
\item Differentiate the function: Find $f'(x)$.
\item Set the derivative equal to zero: $f'(x)=0$.
\item Find the abscissa of the vertex by solving the equation $f'(x)=0$ for $x$ to find the critical $x$ value: $x=k$.
\item Find the ordinate of the vertex by substituting the value of critical value $x=k$ into the function $f(x)$: Evaluate $f(k)$
\end{enumerate}

%-=-=-= EXAMPLE
\begin{example}[id:20150923-152515] \label{20150923-152515}\index{Example!20150923-152515} \hfill \\
Find the vertex of the parabola $y=x^2-2x-6$ using differentiation.
\soln
\solnsteps
1. Differentiate the function.
\begin{align*}
f(x) &=x^2-2x-6 \\
f(x) &=x^2+ \neg 2x + \neg 6 && \text{DOS} \eqref{eq:dos1} \\
\farg{f(x)}' &= \farg{x^2+ \neg 2x + \neg 6}' && \text{SPE} \eqref{eq:spe} \\
f'(x) &= \farg{x^2}'+ \farg{\neg 2x}' + \farg{\neg 6}' && \text{DS} \eqref{eq:ds1} \\
f'(x) &= \farg{x^2}'+ \neg 2 \farg{x}' + \farg{\neg 6}' && \text{DCM} \eqref{eq:dcm1} \\
f'(x) &= 2x + \neg 2 + \farg{\neg 6}' && \text{DPo} \eqref{eq:dpo1} \\
f'(x) &= 2x + \neg 2 + 0 && \text{DC} \eqref{eq:dc1} \\
f'(x) &= 2x + \neg 2 && \text{AId} \eqref{eq:aid2} \\
f'(x) &= 2x - 2 && \text{DOS} \eqref{eq:dos2}
\end{align*}

2 and 3. Set the derivative equal to zero and solve for $x$

\begin{align*}
2x-2 & = 0\\
x &= 1
\end{align*}

4. Find the value of $f(1)$

\begin{align*}
f(x) & = x^2-2x-6 \\
f(1) &= \farg{1}^2-2\farg{1}-6 && \text{SPE} \eqref{eq:spe} \\
f(1) &= -7  &&\text{Evaluate}
\end{align*}

The vertex of this parabola is the point $(1,-7)$
\end{example}


\begin{arule}[Derivative of a Product (DPr)]\index{Derivatives! Derivative of a Product}
\begin{align}
	\left[f(x)g(x) \right]' &= \cRed{f'(x)}g(x) + f(x)\cRed{g'(x)} \label{eq:dpr1} \\
	\dfrac{\dd}{\dx} \left[f(x)g(x) \right] &= \cRed{\dfrac{\dd}{\dx} \left[f(x) \right]}g(x) + f(x)\cRed{\dfrac{\dd}{\dx} \left[g(x)\right]}  \label{eq:dpr2}
\end{align}
\end{arule}

%-=-=-= EXAMPLE
\begin{example}[id:20141209-144203] \label{20141209-144203}\index{Example!20141209-144203} \hfill \\

Differentiate $f(x)=x^2(2x+4)$

\soln

\solnsteps
\begin{align*}
f'(x)&= \farg{x^2(2x+4)}' && \text{SPE} \eqref{eq:spe} \\
f'(x)&= \farg{x^2}'(2x+4)+x^2\farg{2x+4}' && \text{DPr} \eqref{eq:dpr1} \\
f'(x)&= \farg{x^2}'(2x+4)+x^2\farg{2x}+\farg{4}' && \text{DS} \eqref{eq:ds1} \\
f'(x)&= \farg{x^2}'(2x+4)+x^2 \cdot 2 \farg{x}+\farg{4}' && \text{DCM} \eqref{eq:dcm1} \\
f'(x)&= 2x(2x+4)+ x^2\cdot 2 \cdot 1 + \farg{4}' && \text{DPo} \eqref{eq:dpo1} \\
f'(x)&= 2x(2x+4)+ x^2\cdot 2 \cdot 1 + 0 && \text{DC} \eqref{eq:dc1} \\
f'(x)&= 6x^2+8x && \text{simplify} \text{\, goto \,} \, \ref{20141209-145211}\\
\end{align*}
\end{example}


%-=-=-= EXAMPLE
\begin{example}[id:20141209-142321] \label{20141209-142321}\index{Example!20141209-142321} \hfill \\

Differentiate $f(x)=x^2 \cos(x)$

\soln

\solnsteps
\begin{align*}
f'(x)&= \farg{x^2 \cos(x)}' && \text{SPE} \eqref{eq:spe} \\
f'(x)&= \farg{x^2}' \cos(x) + x^2 \farg{\cos (x)}' && \text{DPr} \eqref{eq:dpr1} \\
f'(x)&= 2x \cos (x) + x^2 \farg{\cos (x)}' && \text{DPo} \eqref{eq:dpo1} \\
f'(x)&= 2x \cos (x) + x^2 (-1 \sin(x)) && \text{DCos} \eqref{eq:dcos1} \\
f'(x)&= 2x \cos (x) - x^2 \sin x && \text{OOM} \eqref{eq:oom}
\end{align*}
\end{example}

%-=-=-=-=-=-=-=-=-=-=-=-=-=-=-=-=-=-=-=-=-=-=-=-=
%	SECTION:
%-=-=-=-=-=-=-=-=-=-=-=-=-=-=-=-=-=-=-=-=-=-=-=-=
\section{Derivative of Trigonometric Functions}

\begin{arule}[Derivative of Sine (DSin)]\index{Derivatives! Derivative of Sine}
\begin{align}
	\left[ \sin(x) \right]' &= \cos (x) \label{eq:dsin1} \\
	\dfrac{\dd}{\dx} \left[ \sin(x) \right] & = \cos (x)  \label{eq:dsin2}
\end{align}
\end{arule}


\begin{arule}[Derivative of Cosine (DCos)]\index{Derivatives! Derivative of Cosine}
\begin{align}
	\left[ \cos(x) \right]' &= -\sin (x) \label{eq:dcos1} \\
	\dfrac{\dd}{\dx} \left[ \cos(x) \right] & = -\sin (x)  \label{eq:dcos2}
\end{align}
\end{arule}

%-=-=-= EXAMPLE
\begin{example}[id:20150910-115935] \label{20150910-115935}\index{Example!20150910-115935} \hfill \\

Differentiate $f(x)=\sin(x)\cos(x)$

\soln

\solnsteps
\begin{align*}
f'(x) & = \farg{\sin(x)\cos(x)}' && \text{SPE} \eqref{eq:spe} \\
f'(x) & = \farg{\sin(x)}'\cos(x)+ \sin(x)\farg{\cos(x)}' && \text{DPr} \eqref{eq:dpr1} \\
f'(x) & = \cos (x) \cos (x) + \sin(x)\farg{\cos(x)}' && \text{DSin} \eqref{eq:dsin1} \\
f'(x) &= \cos (x) \cos (x) + \sin(x) (-\sin(x)) && \text{DCos} \eqref{eq:dcos1} \\
f'(x) &= \cos^2(x)-\sin^2(x) && \text{simplify} \text{\, goto \,} \, \ref{} %TODO simplify trigonometric expression
\end{align*}

\end{example}

%-=-=-= EXAMPLE
\begin{example}[id:20141209-151354] \label{20141209-151354}\index{Example!20141209-151354} \hfill \\

Differentiate $f(x)=\sin(x)\sin(x)$

\soln

\solnsteps
\begin{align*}
f'(x) &= \farg{\sin(x)\sin(x)}' && \text{SPE} \eqref{eq:spe} \\
f'(x) &= \farg{\sin(x)}' \sin(x) + \sin (x) \farg{\sin(x)} && \text{DPr} \eqref{eq:dpr1} \\
f'(x) &= \cos(x) \sin(x) + \sin (x) \cos(x) && \text{DSin} \eqref{eq:dsin1} \\
f'(x) &= \cos(x)\sin(x) + \cos(x)\sin(x) && \text{CPM} \eqref{eq:cpm} \\
f'(x) &= 2\cos(x)\sin(x) && \text{OOA} \eqref{eq:ooa} % TODO Simplifying trigonometric expressions
\end{align*}
\end{example}

%-=-=-=-=-=-=-=-=-=-=-=-=-=-=-=-=-=-=-=-=-=-=-=-=
%	SECTION:
%-=-=-=-=-=-=-=-=-=-=-=-=-=-=-=-=-=-=-=-=-=-=-=-=
\section{Derivative of Rational Functions}


\begin{arule}[Derivative of a Quotient (DQ)]\index{Derivatives! Derivative of a Quotient}
\begin{align}
	\left[ \dfrac{f(x)}{g(x)} \right]' &= \dfrac{f'(x)g(x)-f(x)g'(x)}{\left[g(x)\right]^2} \label{eq:drf1} \\
	\dfrac{\dd}{\dx} \left[ \dfrac{f(x)}{g(x)} \right] & = \dfrac{\dfrac{\dd}{\dx}\left[f(x)\right]g(x)-f(x) \dfrac{\dd}{\dx} \left[g(x)\right] }{\left[g(x)\right]^2}  \label{eq:drf2}
\end{align}
\end{arule}



%-=-=-=-=-=-=-=-=-=-=-=-=-=-=-=-=-=-=-=-=-=-=-=-=
%	SECTION:
%-=-=-=-=-=-=-=-=-=-=-=-=-=-=-=-=-=-=-=-=-=-=-=-=
\section{Derivative of Exponential Functions}

%-=-=-=-=-=-=-=-=-=-=-=-=-=-=-=-=-=-=-=-=-=-=-=-=
%	SECTION:
%-=-=-=-=-=-=-=-=-=-=-=-=-=-=-=-=-=-=-=-=-=-=-=-=
\section{Derivative of Logarithmic Functions}

\begin{arule}[Derivative of a Natural Logarithm (DNL)]\index{Derivatives! Derivative of a Natural Logarithm}
\begin{align}
	\left[\ln x \right]' &= \frac{1}{x} \label{eq:dnl1} \\
	\dfrac{\dd}{\dx} \left[\ln x \right] &= \frac{1}{x}  \label{eq:dnl2}
\end{align}
\end{arule}




%-=-=-=-=-=-=-=-=-=-=-=-=-=-=-=-=-=-=-=-=-=-=-=-=
%	SECTION:
%-=-=-=-=-=-=-=-=-=-=-=-=-=-=-=-=-=-=-=-=-=-=-=-=
\section{Derivative of Composite Functions}


\begin{arule}[Derivative of a Composite Function (DComp)]\index{Derivatives! Derivative of a Composite Function}
\begin{align}
	\left[f \left( \cRed{g(x)} \right) \right]' &= \left[\cRed{g(x)} \right]' \left[f\left(\cRed{g(x)}\right)\right]' \label{eq:dcf1} \\
	\dfrac{\dd}{\dx} \left[f \left( \cRed{g(x)} \right) \right] &= \dfrac{\dd}{\dx} \left[\cRed{g(x)} \right] \dfrac{\dd}{\dx} \left[f\left(\cRed{g(x)}\right)\right]  \label{eq:dcf2}
\end{align}
\end{arule}

%-=-=-= EXAMPLE
\begin{example}[id:20141124-203850] \label{20141124-203850}\index{Example!20141124-203850} \hfill \\

Differentiate $y=\ln (3x)$

\soln

\solnsteps

After identifying that $y=\ln (\cRed{3x})$ is a composite function, we let $\cRed{u=3x}$ and thus we get a new function $y=\ln (\cRed{u})$.\\

Using chain rule,
\begin{align*}
\dydx &= \cBlue{\dydu} \cdot \cPurple{\dudx} && \text{DComp} \eqref{eq:dcf2}
\end{align*}

We need to find the factors $\cBlue{\dydu}$ and $\cPurple{\dudx}$.

\begin{alignat*}{2}
		y &= \ln (\cRed{u})				&\quad	 		\cRed{u}		&= 3x \\
	\cBlue{\dydu} &= \dfrac{1}{u}		&\quad 	 \cPurple{\dudx}	 &= 3
\end{alignat*}

\begin{align*}
\dydx &= \cBlue{\dydu} \cdot \cPurple{\dudx} && \text{DComp} \eqref{eq:dcf2} \\
\dydx &= \dfrac{1}{\cRed{u}} \cdot \dudx && \text{DNL} \eqref{eq:dnl1} \\
\dydx &= \dfrac{1}{\cRed{u}} \cdot 3 && \text{DPo} \eqref{eq:dpo1} \\
\dydx &= \dfrac{1}{\cRed{3x}} \cdot 3 \\	% TODO variable subsitution
\dydx &= \dfrac{3}{3x} && \text{OOM} \eqref{eq:oom} \\
\dydx &= \dfrac{1}{x} % TODO reduce fraction
\end{align*}

\qdepend

\qdependlist

example \ref{20141124-205219}-20141124-205219

\end{example}


%-=-=-= EXAMPLE
\begin{example}[id:20141128-160248] \label{20141128-160248}\index{Example!20141128-160248} \hfill \\

Differentiate $y=\ln(3x^2-6x+4)$

\soln

\solnsteps
After identifying that $y=\ln (\cRed{3x^2-6x+4})$ is a composite function, we let \\
$\cRed{u=3x^2-6x+4}$ and thus we get a new function $y=\ln (\cRed{u})$.\\

Using chain rule,
\begin{align*}
\dydx &= \cBlue{\dydu} \cdot \cPurple{\dudx} && \text{DComp} \eqref{eq:dcf2}
\end{align*}

We need to find the factors $\cBlue{\dydu}$ and $\cPurple{\dudx}$.

\begin{alignat*}{3}
		y &= \ln (\cRed{u})				&\quad	 			\cRed{u} 	&= 3x^2-6x+4 	&\quad\\
	\cBlue{\dydu} &= \dfrac{1}{u}		&\quad 	 \cPurple{\dudx}	 &= 6x-6 		&\quad \text{\, goto \,} \, \ref{20141128-151834}
\end{alignat*}

\begin{align*}
\dydx &= \cBlue{\dydu} \cdot \cPurple{\dudx} && \text{DComp} \eqref{eq:dcf2} \\
\dydx &= \dfrac{1}{\cRed{u}} \cdot \dudx && \text{DNL} \eqref{eq:dnl1} \\
\dydx &= \dfrac{1}{\cRed{u}} \cdot (6x-6) && \text{DPo} \eqref{eq:dpo1} \\
\dydx &= \dfrac{1}{\cRed{3x^2-6x+4}} \cdot (6x-6) \\	% TODO variable subsitution
\dydx &= \dfrac{6x-6}{3x^2-6x+4} && \text{OOM} \eqref{eq:oom}
\end{align*}
\end{example}

%-=-=-= EXAMPLE
\begin{example}[id:20141128-155506] \label{20141128-155506}\index{Example!20141128-155506} \hfill \\

Differentiate $y=\ln(\cos x)$

\soln

\solnsteps
After identifying that $y=\ln (\cRed{\cos x})$ is a composite function, we let $\cRed{u=\cos x}$ and thus we get a new function $y=\ln (\cRed{u})$.\\

Using chain rule,
\begin{align*}
\dydx &= \cBlue{\dydu} \cdot \cPurple{\dudx} && \text{DComp} \eqref{eq:dcf2}
\end{align*}

We need to find the factors $\cBlue{\dydu}$ and $\cPurple{\dudx}$.

\begin{alignat*}{2}
		y &= \ln (\cRed{u})				&\quad	 		\cRed{u}		&= \cos x \\
	\cBlue{\dydu} &= \dfrac{1}{u}		&\quad 	 \cPurple{\dudx}	 &= - \sin x
\end{alignat*}

\begin{align*}
\dydx &= \cBlue{\dydu} \cdot \cPurple{\dudx} && \text{DComp} \eqref{eq:dcf2} \\
\dydx &= \dfrac{1}{\cRed{u}} \cdot \dudx && \text{DNL} \eqref{eq:dnl1} \\
\dydx &= \dfrac{1}{\cRed{u}} \cdot - \sin x && \text{DPo} \eqref{eq:dpo1} \\
\dydx &= \dfrac{1}{\cRed{\cos x}} \cdot - \sin x \\	% TODO variable subsitution
\dydx &= \dfrac{-\sin x}{\cos x} && \text{OOM} \eqref{eq:oom} \\
\dydx &= -\tan x % TODO reduce fraction
\end{align*}
\end{example}

%-=-=-= EXAMPLE
\begin{example}[id:20141124-205219] \label{20141124-205219}\index{Example!20141124-205219} \hfill \\

Differentiate $y=\left(x^2-1 \right) \ln(3x)$

\soln

\solnsteps
\begin{align*}
y' &= \farg{x^2-3}' \cdot \ln x + \left(x^2-3\right) \cdot \farg{\ln 3x}' && \text{DPr} \eqref{eq:dpr1} \\
y' &= 2x \cdot \ln x + \left(x^2-3\right) \cdot \farg{\ln 3x}' && \text{differentiate} \text{\, goto \,} \, \ref{20141124-152503} \\
y' &= 2x \cdot \ln x + \left(x^2-3\right) \cdot \dfrac{1}{x}  &&\text{differentiate} \text{\, goto \,} \, \ref{20141124-203850}\\
y' &= 2x \cdot \ln x + \dfrac{x^2-3}{x} && \text{OOM} \eqref{eq:oom} \\
y' &= 2x^2 \ln x +\dfrac{x^2-3}{x} && \text{JTC} \eqref{eq:jtc} \\
y' &= \dfrac{2x^2 \ln x + \left(x^2-3\right)}{x} && \text{OOA} \eqref{eq:ooa}
\end{align*}
\end{example}

\end{document}


%%!TEX root = /Users/markholson/Dropbox/+Projects/LatexFiles/MathNotebook/20150312-163541-rs2.2N-MarksMathNotebook
%-=-=-=-=-=-=-=-=-=-=-=-=-=-=-=-=-=-=-=-=-=-=-=-=
%
%	CHAPTER 
%
%-=-=-=-=-=-=-=-=-=-=-=-=-=-=-=-=-=-=-=-=-=-=-=-=

\chapterimage{chapter_head_2.pdf} % Chapter heading image

\chapter{Integration}

%-=-=-=-=-=-=-=-=-=-=-=-=-=-=-=-=-=-=-=-=-=-=-=-=
%	SECTION: Integration Rules
%-=-=-=-=-=-=-=-=-=-=-=-=-=-=-=-=-=-=-=-=-=-=-=-=

\section{Integration Rules}

\begin{arule}[Integral of a Constant (IC)]\index{Integrals! Integral of a Constant}
\begin{align}
	\displaystyle \int \! c \, \mathrm{d} x &= cx + k \label{eq:ic}
\end{align}
\end{arule}

\begin{arule}[Integral of a Constant Multiple (ICM)]\index{Integrals! Integral of a Constant Multiple}
\begin{align}
	\displaystyle \int \! \alert{c}x \, \mathrm{d} x &= \alert{c} \displaystyle \int \! x \, \mathrm{d} x \label{eq:icm}
\end{align}
\end{arule}

\begin{arule}[Integral of a Power (IPo)]\index{Integrals! Integral of a Power}
\begin{align}
	\displaystyle \int \! x^n \, \mathrm{d} x &=  \dfrac{x^{n+1}}{n+1}  && \text{for} \, x \ne -1 \label{eq:ipo}
\end{align}
\end{arule}

\begin{arule}[Integral of a Sum (IS)]\index{Integrals! Integral of a Sum}
\begin{align}
	\displaystyle \int \! f(x)+g(x) \, \mathrm{d} x &=  \displaystyle \int \! f(x) \, \mathrm{d} x + \displaystyle \int \! g(x) \, \mathrm{d} x  \label{eq:is}
\end{align}
\end{arule}

\begin{example}[id:20141105-172103] \label{20141105-172103} \index{Example!20141105-172103} \hfill \\
	
Integrate $\displaystyle \int \! \dfrac{1}{4}x^2-2x+4 \, \mathrm{d} x$ 

\soln

\solnsteps
\begin{align*}
\displaystyle \int \! \dfrac{1}{4}x^2-2x+4 \, \mathrm{d} x &= \displaystyle \int \! \dfrac{1}{4}x^2 + \neg 2x+4 \, \mathrm{d} x && \text{DOS} \eqref{eq:dos1} \\
&= \displaystyle \int \! \dfrac{1}{4}x^2 \, \mathrm{d} x + \displaystyle \int \! \neg 2x \, \mathrm{d} x + \displaystyle \int \! 4 \, \mathrm{d} x && \text{IS} \eqref{eq:is} \\
&= \dfrac{1}{4}\displaystyle \int \! x^2 \, \mathrm{d} x + \neg 2\displaystyle \int \! x \, \mathrm{d} x + \displaystyle \int \! 4 \, \mathrm{d} x && \text{ICM} \eqref{eq:icm} \\
&= \dfrac{1}{4}\displaystyle \int \! x^2 \, \mathrm{d} x + \neg 2\displaystyle \int \! x \, \mathrm{d} x + 4(x) + c_1 && \text{IC} \eqref{eq:ic} \\
&= \dfrac{1}{4} \left( \dfrac{x^3}{3} \right) + c_2 + \neg 2 \left(\dfrac{x^2}{2} \right) + c_3 + 4(x) + c_1 && \text{IPo} \eqref{eq:ipo} \\
&= \dfrac{x^3}{12} + c_2 - x^2 + c_3 + 4x + c_1 && \text{OOM} \eqref{eq:oom} \\
&= \dfrac{x^3}{12} - x^2 + 4x + \underbrace{c_1 + c_2 + c_3}_{\text{Just a constant}} && \text{CPA} \eqref{eq:cpa} \\
&= \dfrac{x^3}{12} - x^2 + 4x + C && \text{OOA} \eqref{eq:ooa} 
\end{align*}

\qdepend

\qdependlist

example \ref{20141105-144223}-20141105-144223


\end{example}

\begin{theorem}[Fundamental Theorem of Calculus (FThmC)]\index{Theorems!Fundamental Theorem of Calculus}
If $f$ is continuous on the closed interval $[a, b]$ and $F$ is the indefinite integral of $f$ on $[a, b]$, then

\begin{subequations}
\begin{align}
\displaystyle \int_{a}^{b} \! f(x) \, \mathrm{d} x 	&= F(b)-F(a) \label{eq:fthmc}\\
													&= \left[F(x)\right]_{a}^b 
\end{align}
\end{subequations}
\end{theorem}

%-=-=-=-=-=-=-=-=-=-=-=-=-=-=-=-=-=-=-=-=-=-=-=-=
%	SECTION: Volumes of Revolution
%-=-=-=-=-=-=-=-=-=-=-=-=-=-=-=-=-=-=-=-=-=-=-=-=

\section{Volumes of Revolution}

\begin{definition}[Volume of Revolution: Method of Disks]\index{Method of Disks}
Let $f$ be a nonnegative and continuous function on the closed interval $[a,b]$, then the solid of revolution obtained by rotating the curve $f(x)$ about the $x$-axis from $x=a$ to $x=b$ has volume given by

\begin{align}
V & =\pi \displaystyle \int_{a}^{b} \! [f(x)]^2 \, \mathrm{d} x \label{eq:vrd}
\end{align}	
\end{definition}

%-=-=-= EXAMPLE
\begin{example}[id:20141105-144223]\label{20141105-144223} \index{Example!20141105-144223} \hfill \\
	
Find the volume of the solid formed by rotating the region enclosed by $y=2-\dfrac{x}{2}$ and $y=0$ through $2\pi$ about the $x$-axis on the closed interval $[2, 3]$. 

\soln

\solnsteps
\begin{align*}
V & =\pi \displaystyle \int_{a}^{b} \! [f(x)]^2 \, \mathrm{d} x &&\text{VRD}\eqref{eq:vrd} \\
  & =\pi \displaystyle \int_{2}^{3} \! \farg{2-\dfrac{x}{2}}^2 \, \mathrm{d} x && \text{SPE} \eqref{eq:spe} \\
  &= \pi \displaystyle \int_{2}^{3} \! \dfrac{1}{4} x^2 + \neg 2 x + 4 \, \mathrm{d} x && \text{simplify:} \text{\, goto \,} \, \ref{20141105-161225} \\
  &= \pi \left[ \dfrac{x^3}{12} - x^2 + 4x\right]_{2}^{3} && \text{integrate} \text{\, goto \,} \, \ref{20141105-172103} \\
  &= \pi \left[\dfrac{\farg{3}^3}{12} - \farg{3}^2 + 4\farg{3}- \left(\dfrac{\farg{2}^3}{12} - \farg{2}^2 + 4\farg{2} \right) \right] && \text{FThmC} \eqref{eq:fthmc} \\
  &= \pi \left(\dfrac{7}{12}\right)  &&\text{evaluate: example} \, \ref{20141106-083703} \\
  &= \dfrac{7}{12} \pi && \text{CPM} \eqref{eq:cpm}
\end{align*}

% TODO Add tikz plot

\end{example}

%-=-=-= EXAMPLE
\begin{example}[id:20141106-114907] \label{20141106-114907} \index{Example!20141106-114907} \hfill \\
	
Find the volume of the solid formed by rotating the region enclosed by $y=\sqrt{9-x^2}$ and $y=0$ through $2\pi$ about the $x$-axis on the closed interval $[1, 3]$.

\soln

\solnsteps
\begin{align*}
V & = \pi \displaystyle \int_{a}^{b} \! [f(x)]^2 \, \mathrm{d} x &&\text{VRD}\eqref{eq:vrd} \\
  & = \pi \displaystyle \int_{1}^{3} \! \farg{\sqrt{9-x^2}}^2 \, \mathrm{d} x && \text{SPE} \eqref{eq:spe} \\
  & = \pi \displaystyle \int_{1}^{3} \! 9-x^2 \, \mathrm{d} x && \text{simplify:} \text{\, goto \,} \, \ref{20141107-121834}\\
  & = \pi \left[9x-\dfrac{x^3}{3} \right]_{1}^{3}  &&\text{integrate} \\ % TODO Create integration example
  & = \pi \left[9\farg{3}-\dfrac{\farg{3}^3}{3} - \left(9\farg{1}-\dfrac{\farg{1}^3}{3} \right) \right] && \text{FThmC} \eqref{eq:fthmc} \\
  & = \pi \left(\dfrac{28}{3} \right)  &&\text{evaluate} \\
  & = \dfrac{28}{3}\pi && \text{CPM} \eqref{eq:cpm}  
\end{align*}

% TODO Add tikz plot

\end{example}

%-=-=-= EXAMPLE
\begin{example}[id:20141106-122528] \label{20141106-122528} \index{Example!20141106-122528} \hfill \\
	
Find the volume of the solid formed by rotating the region enclosed by $y=\sqrt{\sin x}$ and $y=0$ through $2\pi$ about the $x$-axis on the closed interval $\left[\frac{\pi}{2}, \pi \right]$.

\soln

\solnsteps
\begin{align*}
V & = \pi \displaystyle \int_{a}^{b} \! [f(x)]^2 \, \mathrm{d} x &&\text{VRD}\eqref{eq:vrd} \\
  & = \pi \displaystyle \int_{\pi/2}^{\pi} \! \farg{\sqrt{\sin} \theta}^2 \, \mathrm{d} \theta && \text{SPE} \eqref{eq:spe} \\
  & = \pi \displaystyle \int_{\pi/2}^{\pi} \! \sin \theta  \, \mathrm{d} \theta && \text{PoPo} \eqref{eq:popo1} \\
  & = \pi \left[ -1 \cos \theta \right]_{\pi/2}^{\pi}  &&\text{integrate} \\
  & = \pi \left[-1\cos \farg{\pi}- \left(-1\cos \farg{\dfrac{\pi}{2}} \right) \right] && \text{FThmC} \eqref{eq:fthmc} \\
  & = \pi (1)  &&\text{evaluate} \\
  & = \pi && \text{MId} \eqref{eq:mid2}  
\end{align*}

% TODO Add tikz plot

\end{example}

%-=-=-= EXAMPLE
\begin{example}[id:20141108-083108] \label{20141108-083108} \index{Example!20141108-083108} \hfill \\

Find the volume of the solid formed by rotating the region enclosed by $y=\sqrt{2-x^2}$ and $y=0$ through $2\pi$ about the $x$-axis.

\soln
\solnsteps


Since no closed interval is stated, we will need to find where the curves $y=\sqrt{2-x^2}$ and $y=0$ intersect to determine the closed interval.  Solving this system of equations (see example \ref{20141107-131748}), we find that $x=\pm \sqrt{2}$.

\begin{align*}
V & = \pi \displaystyle \int_{a}^{b} \! [f(x)]^2 \, \mathrm{d} x &&\text{VRD}\eqref{eq:vrd} \\
&= \pi \displaystyle \int_{-\sqrt{2}}^{\sqrt{2}} \! \farg{\sqrt{2-x^2}}^2 \, \mathrm{d} x \\ % TODO Justify step
&= \pi \displaystyle \int_{-\sqrt{2}}^{\sqrt{2}} \! 2-x^2 \, \mathrm{d} x && \text{PoPo} \eqref{eq:popo1} \\
&= \pi \left[2x-\dfrac{x^3}{3} \right]_{-\sqrt{2}}^{\sqrt{2}}  &&\text{integrate} \\ % TODO add example
&= \pi \left[2\farg{\sqrt{2}}-\dfrac{\farg{\sqrt{2}}^3}{3}-\left(2\farg{-\sqrt{2}}-\dfrac{\farg{-\sqrt{2}}^3}{3} \right) \right] && \text{FThmC} \eqref{eq:fthmc} \\
&= \pi \left(\dfrac{8}{2} \sqrt{2}\right)  &&\text{simplify} \text{\, goto \,} \, \ref{20141107-121834}\\
&= \dfrac{8}{2} \pi \sqrt{2} && \text{CPM} \eqref{eq:cpm} 
\end{align*}

% TODO Add tikz plot

\end{example}

%-=-=-= EXAMPLE
\begin{example}[id:20141125-084541] \label{20141125-084541}\index{Example!20141125-084541} \hfill \\

Find the volume of the solid formed by rotating the region enclosed by $y=\sqrt{2-x^2}$ and $y=0$ through $2\pi$ about the $x$-axis.

\soln
\solnsteps

Since no closed interval is stated, we will need to find where the curves $y=4-x^2$ and $y=0$ intersect to determine the closed interval.  Solving this system of equations (see the similar example \ref{20141107-131748}), we find that $x=\pm 2$. \\


Since the parabolar $y=4-x^2$ is symmetric about the $y$-axis, we can find the volume on the closed interval [0,2] and multiply this by a factor of two to include the volume on the closed interval [-2,0].
 
\soln

\solnsteps
\begin{align*}
V & = 2 \cdot \pi \displaystyle \int_{a}^{b} \! [f(x)]^2 \, \mathrm{d} x &&\text{VRD}\eqref{eq:vrd} \\
&= 2 \cdot \pi \displaystyle \int_{0}^{2} \! \farg{4-x^2}^2 \, \mathrm{d} x \\ % TODO Justify step
&= 2 \cdot \pi \displaystyle \int_{0}^{2} \! x^4-8x^2+16 \, \mathrm{d} x  \\ % TODO Justify step
&= 2 \cdot \pi \left[\dfrac{x^5}{5}-\dfrac{8x^3}{3}+16x \right]_{0}^{2} \\ % TODO Justify step
&= 2 \cdot \pi \left( \dfrac{\farg{2}^5}{5}-\dfrac{8\farg{2}^3}{3}+16\farg{2} \right) && \text{FThmC} \eqref{eq:fthmc} \\
&= \dfrac{512 \pi}{15}
\end{align*}
\end{example}

%%-=-=-=-=-=-=-=-=-=-=-=-=-=-=-=-=-=-=-=-=-=-=-=-=
%
%	CHAPTER: 
%
%-=-=-=-=-=-=-=-=-=-=-=-=-=-=-=-=-=-=-=-=-=-=-=-=
\chapter{GY2011}
\newpage
\pagestyle{empty}
\begin{landscape}
\section{GY2011 Local Knowledge Requirement Summaries}
\label{chap:LocalKnowledgeRequirements}


\newcolumntype{A}{ >{\centering\arraybackslash} m{5cm} }
\newcolumntype{B}{ >{\centering\arraybackslash} m{6cm} }
\newcolumntype{C}{ >{\centering\arraybackslash} m{6cm} }
\newcolumntype{D}{ >{\centering\arraybackslash} m{6cm} }
\begin{longtable}{ | A | B | C | D |}
% START HEADER
\hline 
\rowcolor[gray]{.9}   \textbf{Category} &
    \textbf{E Local Summary} &
    \textbf{C Local Summary} &
    \textbf{A Local Summary}\\
  \hline  
\endhead  
% END HEADER


%CONCEPT
\textbf{Concept} &
\alert{Understands} concepts, their purpose and connections between them, illustrated by \alert{some} examples and different representations &
\alert{Thoroughly understands} concepts, their purpose and connections between them, illustrated by \alert{some} examples and different representations &
\alert{Thoroughly understands} concepts, their purpose and connections between them, illustrated by \alert{several} examples and different representations
\\

\hline

%PROCEDURE
\textbf{Procedure} &
Carries out one-step procedures, \alert{some} alternative solution methods &
Carries out \alert{multi-step} procedures&
\alert{Fluently interchanges} solution methods
\\

\hline

%PROBLEM SOLVING
\textbf{Problem Solving} &
Solves problems in \alert{familiar} situations, with \alert{simple} interpretations/formulations &
Solves, formulates and analyses problems with \alert{advanced} interpretations/formulations&
Solves, formulates and analyses \alert{complex problems effectively}
\\

\hline

%MODELING
\textbf{Modeling} &
Uses given models &
\alert{Chooses and applies} appropriate models&
\alert{Adapts} mathematical models and \alert{provides alternative} solution strategies
\\

\hline

%REASONING AND EVALUATION
\textbf{Reasoning \& Evaluation} &
\alert{Simple} evaluation and reasoning, differentiates between guesses and well-founded statements &
\alert{Well-founded} reasoning with \alert{nuanced} judgement and evaluation&
\alert{Well-founded} and \alert{nuanced} reasoning, judgement and evaluation. Finds \alert{general rules and relationships} and expresses these using \alert{algebra and symbols}
\\

\hline

%COMMUNICATIONS
\textbf{Communication} &
Can communicate with \alert{some} certainty. Uses symbols, terms and conventions &
Can communicate with \alert{some} certainty. Uses symbols, terms and conventions \alert{according to purpose and context}&
Can communicate with certainty. Symbols, terms and conventions are \alert{well-suited} to purpose and context
\\

\hline

%CONTEXT AND RELEVANCE
\textbf{Context \& Relevance} &
\alert{Some} contextualization through examples &
\alert{Relevant} contextualization through examples &
\alert{Well-founded} and \alert{nuanced} contextualization through examples 
\\

\hline
\caption{Table caption}
\end{longtable}


\end{landscape}
%%------------------------------------------------

\section{Citation}\index{Citation}

This statement requires citation \cite{book_key}; this one is more specific \cite[122]{article_key}.

%------------------------------------------------

\section{Lists}\index{Lists}

Lists are useful to present information in a concise and/or ordered way\footnote{Footnote example...}.

\subsection{Numbered List}\index{Lists!Numbered List}

\begin{enumerate}
\item The first item
\item The second item
\item The third item
\end{enumerate}

\subsection{Bullet Points}\index{Lists!Bullet Points}

\begin{itemize}
\item The first item
\item The second item
\item The third item
\end{itemize}

\subsection{Descriptions and Definitions}\index{Lists!Descriptions and Definitions}

\begin{description}
\item[Name] Description
\item[Word] Definition
\item[Comment] Elaboration
\end{description}

%----------------------------------------------------------------------------------------
%	CHAPTER 2
%----------------------------------------------------------------------------------------

\chapter{In-text Elements}

\lipsum[1-7] % Dummy text

\section{Theorems}\index{Theorems}

This is an example of theorems.

\subsection{Several equations}\index{Theorems!Several Equations}
This is a theorem consisting of several equations.

\begin{theorem}[Name of the theorem]
In $E=\mathbb{R}^n$ all norms are equivalent. It has the properties:
\begin{align}
& \big| ||\mathbf{x}|| - ||\mathbf{y}|| \big|\leq || \mathbf{x}- \mathbf{y}||\\
&  ||\sum_{i=1}^n\mathbf{x}_i||\leq \sum_{i=1}^n||\mathbf{x}_i||\quad\text{where $n$ is a finite integer}
\end{align}
\end{theorem}

\subsection{Single Line}\index{Theorems!Single Line}
This is a theorem consisting of just one line.

\begin{theorem}
A set $\mathcal{D}(G)$ in dense in $L^2(G)$, $|\cdot|_0$. 
\end{theorem}

%------------------------------------------------

\section{Definitions}\index{Definitions}

This is an example of a definition. A definition could be mathematical or it could define a concept.

\begin{definition}[Definition name]
Given a vector space $E$, a norm on $E$ is an application, denoted $||\cdot||$, $E$ in $\mathbb{R}^+=[0,+\infty[$ such that:
\begin{align}
& ||\mathbf{x}||=0\ \Rightarrow\ \mathbf{x}=\mathbf{0}\\
& ||\lambda \mathbf{x}||=|\lambda|\cdot ||\mathbf{x}||\\
& ||\mathbf{x}+\mathbf{y}||\leq ||\mathbf{x}||+||\mathbf{y}||
\end{align}
\end{definition}

%------------------------------------------------

\section{Notations}\index{Notations}

\begin{notation}
Given an open subset $G$ of $\mathbb{R}^n$, the set of functions $\varphi$ are:
\begin{enumerate}
\item Bounded support $G$;
\item Infinitely differentiable;
\end{enumerate}
a vector space is denoted by $\mathcal{D}(G)$. 
\end{notation}

%------------------------------------------------

\section{Remarks}\index{Remarks}

This is an example of a remark.

\begin{remark}
The concepts presented here are now in conventional employment in mathematics. Vector spaces are taken over the field $\mathbb{K}=\mathbb{R}$, however, established properties are easily extended to $\mathbb{K}=\mathbb{C}$.
\end{remark}

%------------------------------------------------

\section{Corollaries}\index{Corollaries}

This is an example of a corollary.

\begin{corollary}[Corollary name]
The concepts presented here are now in conventional employment in mathematics. Vector spaces are taken over the field $\mathbb{K}=\mathbb{R}$, however, established properties are easily extended to $\mathbb{K}=\mathbb{C}$.
\end{corollary}

%------------------------------------------------

\section{Propositions}\index{Propositions}

This is an example of propositions.

\subsection{Several equations}\index{Propositions!Several Equations}

\begin{proposition}[Proposition name]
It has the properties:
\begin{align}
& \big| ||\mathbf{x}|| - ||\mathbf{y}|| \big|\leq || \mathbf{x}- \mathbf{y}||\\
&  ||\sum_{i=1}^n\mathbf{x}_i||\leq \sum_{i=1}^n||\mathbf{x}_i||\quad\text{where $n$ is a finite integer}
\end{align}
\end{proposition}

\subsection{Single Line}\index{Propositions!Single Line}

\begin{proposition} 
Let $f,g\in L^2(G)$; if $\forall \varphi\in\mathcal{D}(G)$, $(f,\varphi)_0=(g,\varphi)_0$ then $f = g$. 
\end{proposition}

%------------------------------------------------

\section{Examples}\index{Examples}

This is an example of examples.

\subsection{Equation and Text}\index{Examples!Equation and Text}

\begin{example}
Let $G=\{x\in\mathbb{R}^2:|x|<3\}$ and denoted by: $x^0=(1,1)$; consider the function:
\begin{equation}
f(x)=\left\{\begin{aligned} & \mathrm{e}^{|x|} & & \text{si $|x-x^0|\leq 1/2$}\\
& 0 & & \text{si $|x-x^0|> 1/2$}\end{aligned}\right.
\end{equation}
The function $f$ has bounded support, we can take $A=\{x\in\mathbb{R}^2:|x-x^0|\leq 1/2+\epsilon\}$ for all $\epsilon\in\intoo{0}{5/2-\sqrt{2}}$.
\end{example}

\subsection{Paragraph of Text}\index{Examples!Paragraph of Text}

\begin{example}[Example name]
\lipsum[2]
\end{example}

%------------------------------------------------

\section{Exercises}\index{Exercises}

This is an example of an exercise.

\begin{exercise}
Simplify the following expressions.
\begin{multicols}{3}
\begin{enumerate}
\item $5x+4$
\columnbreak
\item $5x+4$
\columnbreak
\item $5x+4$
\end{enumerate}
\end{multicols}
\end{exercise}

\begin{exercise}
This is a good place to ask a question to test learning progress or further cement ideas into students' minds.
\end{exercise}

\begin{exercise}
This is a good place to ask a question to test learning progress or further cement ideas into students' minds.
\end{exercise}

%------------------------------------------------

\section{Problems}\index{Problems}

\begin{problem}
What is the average airspeed velocity of an unladen swallow?
\end{problem}

%------------------------------------------------

\section{Vocabulary}\index{Vocabulary}

Define a word to improve a students' vocabulary.

\begin{vocabulary}[Word]
Definition of word.
\end{vocabulary}

%----------------------------------------------------------------------------------------
%	CHAPTER 3
%----------------------------------------------------------------------------------------

\chapterimage{chapter_head_1.pdf} % Chapter heading image

\chapter{Presenting Information}

\section{Table}\index{Table}

\begin{table}[h]
\centering
\begin{tabular}{l l l}
\toprule
\textbf{Treatments} & \textbf{Response 1} & \textbf{Response 2}\\
\midrule
Treatment 1 & 0.0003262 & 0.562 \\
Treatment 2 & 0.0015681 & 0.910 \\
Treatment 3 & 0.0009271 & 0.296 \\
\bottomrule
\end{tabular}
\caption{Table caption}
\end{table}

%------------------------------------------------

\section{Figure}\index{Figure}

\begin{figure}[h]
\centering\includegraphics[scale=0.5]{placeholder}
\caption{Figure caption}
\end{figure}

%-=-=-=-=-=-=-=-=-=-=-=-=-=-=-=-=-=-=-=-=-=-=-=-=
%	BIBLIOGRAPHY
%-=-=-=-=-=-=-=-=-=-=-=-=-=-=-=-=-=-=-=-=-=-=-=-=

\chapter*{Bibliography}
\addcontentsline{toc}{chapter}{\textcolor{ocre}{Bibliography}}
\section*{Books}
\addcontentsline{toc}{chapter}{\textcolor{ocre}{Bibliography}}
\section*{Website}
\addcontentsline{toc}{section}{Websites}
\printbibliography[heading=bibempty,type=misc]
\section*{Articles}
\addcontentsline{toc}{section}{Articles}
\printbibliography[heading=bibempty,type=article]

%-=-=-=-=-=-=-=-=-=-=-=-=-=-=-=-=-=-=-=-=-=-=-=-=
%	INDEX
%-=-=-=-=-=-=-=-=-=-=-=-=-=-=-=-=-=-=-=-=-=-=-=-=

\cleardoublepage
\phantomsection
\setlength{\columnsep}{0.75cm}
\addcontentsline{toc}{chapter}{\textcolor{ocre}{Index}}
\printindex

\end{document}
